\documentclass[11pt]{article}

\usepackage[T1]{fontenc}
\usepackage[utf8]{inputenc}
\usepackage[english]{babel}
\usepackage[margin=1in]{geometry}
\usepackage{lmodern}
\usepackage{microtype}
\usepackage{mathtools}
\usepackage{amsmath,amssymb,amsthm}
\usepackage{array}
\usepackage{booktabs}
\usepackage{tabularx}
\usepackage{longtable}
\usepackage{float}
\usepackage{graphicx}
\usepackage{pdflscape}
\usepackage{caption}
\usepackage{enumitem}
\usepackage{fancyhdr}
\usepackage[numbers]{natbib}
\usepackage{doi}
\usepackage{xurl}
\usepackage{hyperref}
\usepackage{cleveref}
\usepackage{orcidlink}
\usepackage{titling}



\newcommand{\SixBirds}{Six Birds}
\newcommand{\QMgr}{QM--GR}
\newcommand{\Deltafact}{\Delta_{\mathrm{fact}}}
\newcommand{\RecordStable}{\mathrm{RS}}
\newcommand{\PredictionOne}{Prediction~1}
\newcommand{\PredictionTwo}{Prediction~2}
\newcommand{\PredictionThree}{Prediction~3}
 
\DeclareUnicodeCharacter{00A7}{\S}
\DeclareUnicodeCharacter{039B}{\ensuremath{\Lambda}}
\DeclareUnicodeCharacter{03B4}{\ensuremath{\delta}}
\DeclareUnicodeCharacter{2013}{--}
\DeclareUnicodeCharacter{2192}{\ensuremath{\to}}
\DeclareUnicodeCharacter{2194}{\ensuremath{\leftrightarrow}}
\DeclareUnicodeCharacter{21D2}{\ensuremath{\Rightarrow}}

\newtheorem{theorem}{Theorem}[section]
\newtheorem{lemma}[theorem]{Lemma}
\newtheorem{proposition}[theorem]{Proposition}
\newtheorem{corollary}[theorem]{Corollary}
\theoremstyle{definition}
\newtheorem{definition}[theorem]{Definition}
\newtheorem{remark}[theorem]{Remark}
\providecommand{\tightlist}{\setlength{\itemsep}{0pt}\setlength{\parskip}{0pt}}
\urlstyle{same}
\captionsetup{
  font=small,
  labelfont=bf,
  labelsep=period
}
\captionsetup[figure]{skip=6pt}
\crefname{theorem}{theorem}{theorems}
\Crefname{theorem}{Theorem}{Theorems}
\crefname{lemma}{lemma}{lemmas}
\Crefname{lemma}{Lemma}{Lemmas}
\crefname{proposition}{proposition}{propositions}
\Crefname{proposition}{Proposition}{Propositions}
\crefname{corollary}{corollary}{corollaries}
\Crefname{corollary}{Corollary}{Corollaries}
\crefname{definition}{definition}{definitions}
\Crefname{definition}{Definition}{Definitions}
\crefname{remark}{remark}{remarks}
\Crefname{remark}{Remark}{Remarks}
\setlist[enumerate]{leftmargin=*, itemsep=2pt, topsep=4pt}
\setlength{\droptitle}{-3.2em}
\pretitle{\begin{center}\LARGE\bfseries}
\posttitle{\par\end{center}\vspace{0.5em}}
\preauthor{\begin{center}\normalsize}
\postauthor{\par\end{center}\vspace{-0.4em}}
\predate{\begin{center}\small}
\postdate{\par\end{center}\vspace{0.6em}}
\setlength{\emergencystretch}{2em}
\clubpenalty=10000
\widowpenalty=10000
\displaywidowpenalty=10000
\interfootnotelinepenalty=10000
\allowdisplaybreaks[2]
\fancypagestyle{firstpage}{\fancyhf{}
  \fancyfoot[C]{\scriptsize Automorph Inc., Wilmington, DE, USA \quad \texttt{ioannis@automorph.io} \quad \textcopyright\ 2026 \quad CC-BY 4.0 \quad \doi{10.5281/zenodo.20713213}}
  \renewcommand{\headrulewidth}{0pt}
  \renewcommand{\footrulewidth}{0pt}
}

\title{To Kill Three Stones with Six Birds:\\ A Common Grammar for the SM, QM, and GR}
\author{Ioannis Tsiokos\,\orcidlink{0009-0009-7659-5964}}
\date{\today}
\hypersetup{
  pdftitle={To Kill Three Stones with Six Birds: A Common Grammar for the SM, QM, and GR},
  pdfauthor={Ioannis Tsiokos},
  pdfkeywords={Six Birds Theory; emergence calculus; layer-agnostic structural laws; foundations of physics; Standard Model gauge selection; quantum gravity; holographic entanglement; gravitationally-induced entanglement; grand unification; falsifiable predictions},
  hidelinks
}

\begin{document}

\maketitle
\thispagestyle{firstpage}

\begin{abstract}
A standing question in foundations is whether the deep problems of distinct domains --- \emph{why} the Standard Model has
the structure it does, and \emph{why} quantum mechanics and general relativity fail to combine --- share any common
explanatory machinery. We report that one frozen, \textbf{layer-agnostic structural calculus} (the Six Birds emergence
calculus) --- a catalog of substrate-independent laws governing how one theory emerges as a quotient/closure of another
--- produces foundational-grade results on both, and yields \textbf{three forced, falsifiable, mainstream-contradicting
predictions}.

Applied under an anti-contamination guard on finite, audited, reproducible toy carriers, the calculus gives the two
opposite result-genres its own laws anticipate. For the Standard Model, a \emph{demarcation} separating the structurally
forced from the irreducibly observed-input, recovering non-circularly the grand-unified embedding ratios
\(\sin^2\theta_W = 3/8\) and \(k_Y = 5/3\) and unifying proton stability, monopole absence, and color separation as one
object. For QM--GR, a \emph{unification} that recognizes Ryu--Takayanagi holographic entanglement as a case of one
ledger/shadow-price duality.

The three predictions each contradict a prevailing program. \textbf{(i)} Boundary entanglement does \emph{not} determine
bulk geometry: the entanglement\(\to\)geometry map carries a robust kernel, so strong ``it-from-qubit'' fails.
\textbf{(ii)} Gravity does \emph{not} mediate entanglement: its channel is forced to be classically-indexed, so BMV-class
experiments are predicted \textbf{null}, testable on tabletop timescales within the decade. \textbf{(iii)} Record-stability
forbids proton decay and magnetic monopoles, against grand unification. A single record/memory layer ties the predictions
at the grammar level without identifying the two physical tracks.

Every result is computed under adversarial self-correction against explicit can-fail controls, and two construction results
were stress-tested in adversarial review cycles and settled. We claim \textbf{no frame-transfer}: the results are conditional, toy-level structural facts
(\(n=2\) on universality), offered as a proof of principle and placed before the community to test.
\end{abstract}

\noindent\textbf{Keywords.} Six Birds Theory; emergence calculus; layer-agnostic structural laws; foundations of physics; Standard Model gauge selection; quantum gravity; holographic entanglement; gravitationally-induced entanglement; grand unification; falsifiable predictions.

\hypertarget{introduction}{\section{Introduction}\label{introduction}}

Physics carries two foreclosures that have resisted a century of attack from within their own domains.

The first is \emph{selection}: why \textbf{this} Standard Model? The gauge algebra \(\mathfrak{su}(2)\oplus\mathfrak{su}(3)\oplus \mathfrak{u}(1)\)~\citep{Weinberg1967}, the hypercharge pattern, three generations, the electroweak scale, this vacuum --- none of it is forced
by the internal consistency of quantum field theory, and the program that tried to derive it from above (grand
unification~\citep{GeorgiGlashow1974}, then the landscape) has not converged. The second is \emph{composition}: why do quantum mechanics and general
relativity --- each spectacularly confirmed in its own regime --- fail to combine? Quantizing gravity perturbatively fails
(non-renormalizability~\citep{tHooftVeltman1974,GoroffSagnotti1985}); fifty years of candidate completions have produced no confirmed quantum theory of gravity.

These two problems are, on their face, structurally unrelated. One asks why a \emph{contingent-looking} structure was
selected; the other asks why two \emph{law-like} structures do not compose. They are studied by different communities with
different tools, and no standard framework treats them as instances of anything common.

\textbf{The claim of this paper} is that they \emph{are} instances of something common, and that the common machinery is not a
new dynamical theory but a \textbf{layer-agnostic structural calculus}: a fixed, finite catalog of laws about how one
description of the world emerges from another --- as a quotient that forgets detail, audited so that nothing is smuggled.
Pointed at the two problems under an explicit anti-contamination protocol (Section 3), the same frozen calculus
produces, on finite reproducible carriers:

\begin{enumerate}
\def\labelenumi{\arabic{enumi}.}
\item
  \textbf{For the Standard Model --- a demarcation} (Section 4). A proof-grade separation of the structurally forced from the
  irreducibly contingent: an unconditional exclusion of \(\sim 99.3\%\) of genuinely-chiral gauge structures; a
  conditional selection of \(\mathfrak{su}(2)\oplus\mathfrak{su}(3)\oplus\mathfrak{u}(1)\) on one declared, independently
  grounded condition; non-circular recovery of the grand-unified embedding ratios \(\sin^2\theta_W=3/8\) and \(k_Y=5/3\);
  and theorems proving the calculus is \emph{blind} to the generation count and the detailed matter content --- i.e., a proof
  that those are observed-input, not derivable.
\item
  \textbf{For QM--GR --- a unification} (Section 5). The non-composition is located as a \emph{route-mismatch} (two idempotent
  completions, ``quantize'' and ``curve'', that provably fail to commute), and resolved by a \emph{fork}: QM and GR are sibling
  read-outs of one common refinement --- the directed reduction of either to the other provably fails, while the
  symmetric parent closes. Within this structure, the Ryu--Takayanagi relation~\citep{RyuTakayanagi2006} between boundary entanglement and bulk
  area is \emph{recognized} --- not assumed --- as a special case of a single ledger/shadow-price duality (max-flow/min-cut as
  linear-programming duality), a recovery confirmed by a separate adversarial review pass.
\item
  \textbf{Three falsifiable, mainstream-contradicting predictions} (Section 6) --- the paper's principal results:

  \begin{quote}
  \textbf{Prediction 1 (contra strong it-from-qubit).} \emph{Boundary entanglement does not determine bulk geometry.} The
  map from the complete boundary-entanglement fingerprint to the bulk geometry has a robust, finite-dimensional
  kernel: physically distinct geometries share identical entanglement data. Complete entanglement-based bulk
  reconstruction is therefore impossible. \textbf{Falsifier:} a non-trivial holographic/tensor-network geometry whose
  entanglement\(\to\)geometry map is injective.
  \end{quote}

  \begin{quote}
  \textbf{Prediction 2 (contra quantized-mediator gravity; near-term).} \emph{Gravity does not mediate entanglement.} The
  same fork forces the gravitational channel to be a classically-indexed \emph{record} channel --- it reads and records the
  geometry-visible configuration rather than coupling coherently to the phase mode. Two masses coupled through it
  \textbf{decohere without becoming entangled}. \textbf{Falsifier:} a confirmed BMV-type observation of gravitationally-induced
  entanglement between masses (tabletop, plausibly within the decade) --- the nearest-term of the three.
  \end{quote}

  \begin{quote}
  \textbf{Prediction 3 (contra grand unification).} \emph{Record-stability forbids proton decay and magnetic monopoles.} The
  capacity of a universe to bear stable records is structurally tied to baryon-number conservation: any
  record-bearing structure lies in the ``clean'' branch in which the \(X/Y\)-type coset that would mediate proton decay
  and monopoles is absent. \textbf{Falsifier:} an observed proton-decay event (e.g., at Hyper-Kamiokande) or a confirmed
  magnetic monopole --- our universe manifestly bears records, so either observation breaks the predicted link.
  \end{quote}

  In each case the \emph{distinctive} content is structural and contrarian: prevailing programs assert, respectively, that
  entanglement \emph{builds} geometry~\citep{VanRaamsdonk2010,MaldacenaSusskind2013}, that gravity \emph{will} entangle (BMV-positive)~\citep{Bose2017,MarlettoVedral2017}, and that protons \emph{do} decay~\citep{Langacker1981}. The
  calculus forces the opposite in all three, and attaches a falsification condition to each. The three
  are not independent stipulations: a single record/memory layer ties them at the grammar level (it grounds the gauge selection,
  forbids the decay, and classicalizes the gravitational channel) --- a cross-domain structural tie examined in Section 8.
\item
  \textbf{A universality datum} (Sections 7--8). The same defect/obstruction machinery also classifies Bell nonlocality,
  black-hole information loss, and the cosmological constant --- three further deep puzzles --- each as a computed
  forbidden-rule or demarcation. And across the two headline problems, the calculus produced \emph{opposite genres of
  result} (a demarcation; a unification) exactly where its own laws said it would, with that differential expectation
  on record before the constructions were run.
\end{enumerate}

\textbf{What this paper is not.} We do not solve quantum gravity; we do not derive any measured low-energy constant; we make
no claim of frame-transfer from the finite carriers to nature. Every landing in this paper carries an explicit grade
(Section 2.3) and an explicit falsification condition (Section 9, Table T1). Stated plainly, the contribution
is: \emph{a single frozen structural grammar, applied across physics' two deepest and most unrelated problems, produces
foundational-grade structure on both and yields three falsifiable predictions} --- the predictions of Section 6 are
where it is most open to refutation.

\textbf{Reader's map.} Section 2 is a self-contained primer on the calculus (no prior knowledge assumed). Section 3 states
the construction-and-audit protocol, including the adversarial self-correction record. Sections 4--5 present the two track results; Section 6 the predictions; Section 7 the
breadth results; Section 8 the cross-track thesis; Section 9 the consolidated scope/limits and the claim-by-claim
falsifiability table. Appendices give the formal calculus (A), full reproducibility data (B), the rejected-attempts
audit trail (C), the adversarial-review record (D), and notation (E).
 \hypertarget{the-six-birds-emergence-calculus-in-brief}{\section{The Six Birds emergence calculus, in brief}\label{the-six-birds-emergence-calculus-in-brief}}

This section is a self-contained primer; no prior exposure to the framework is assumed. The formal development lives in
the foundations papers (Foundations II--IV~\citep{Tsiokos_FoundationsII,Tsiokos_FoundationsIII,Tsiokos_FoundationsIV}; \emph{Foundations of Emergence Calculus}~\citep{Tsiokos_FoEC}; \emph{Why Mathematics Even Works}~\citep{Tsiokos_WhyMath}) and is
summarized operationally in Appendix A. Here we give the minimum a working physicist needs to read Sections 4--8.

\hypertarget{layers-quotients-and-descent}{\subsection{Layers, quotients, and descent}\label{layers-quotients-and-descent}}

The organizing picture is that physical description is \textbf{layered}~\citep{Anderson1972}: thermodynamics over statistical mechanics, hadrons
over quarks, classical fields over quantum states, geometry over whatever underlies it. The calculus models a layer as
a \textbf{quotient of a carrier}:

\begin{itemize}
\tightlist
\item
  a \textbf{carrier} \(H\) --- a set of finer-grained configurations (microstates, field configurations, tensor data);
\item
  an \textbf{access quotient} \(q : H \to Q\) --- the surjection that forgets what the layer cannot see (two configurations are
  layer-identical iff \(q\) maps them together);
\item
  \textbf{readouts} \(t : H \to T\) --- quantities one wants to evaluate.
\end{itemize}

The central, endlessly reused question is whether a readout \textbf{descends} to the layer:

\[
t \text{ descends through } q
\;\iff\;
\forall h,h':\; q(h)=q(h') \Rightarrow t(h)=t(h')
\;\iff\;
\exists\, \bar t : Q \to T,\;\; \bar t \circ q = t .
\]

When descent fails, the failure is recorded as a finite \textbf{obstruction set}

\[
\mathcal O_q(t) \;=\; \{(h,h') : q(h)=q(h') \ \wedge\ t(h)\neq t(h')\},
\]

and the pair \((h,h')\) is a \emph{witness}: two configurations the layer cannot distinguish that nonetheless differ in \(t\).
Everything in this paper --- selection results, no-go theorems, and the three falsifiable predictions --- is ultimately a
statement about such obstruction sets being empty or non-empty, computed exactly on finite carriers.

A second reusable object is the \textbf{factorization defect} between two maps \(\pi_0,\pi_1\) on a common carrier:

\[
\Delta_{\mathrm{fact}}(\pi_0,\pi_1)
\;=\;
\{(s,s') : \pi_0(s)=\pi_0(s') \ \wedge\ \pi_1(s)\neq\pi_1(s')\},
\qquad
\Delta_{\mathrm{fact}} \neq \varnothing \iff \pi_1 \not\preceq \pi_0,
\]

i.e., \(\pi_1\) fails to factor through \(\pi_0\) exactly when the defect is non-empty (a theorem-grade equivalence;
Appendix A). In Section 4 this single object --- applied to gauge structures --- \emph{is} the \(X/Y\) coset that unifies
proton-stability, monopole-absence, and color-separation; in Section 5 it is the witness that GR is not a quotient of QM.

\hypertarget{the-six-primitive-roles}{\subsection{The six primitive roles}\label{the-six-primitive-roles}}

The calculus fixes a finite alphabet of six \textbf{primitive roles},

\[
\mathbb P = \{P_1,\dots,P_6\},
\]

naming the recurring ways layers relate (Foundations III). Operationally:

\begin{longtable}[]{@{}llll@{}}
\toprule
\begin{minipage}[b]{0.11\columnwidth}\raggedright
primitive\strut
\end{minipage} & \begin{minipage}[b]{0.18\columnwidth}\raggedright
role\strut
\end{minipage} & \begin{minipage}[b]{0.35\columnwidth}\raggedright
one-line gloss\strut
\end{minipage} & \begin{minipage}[b]{0.26\columnwidth}\raggedright
where it appears in this paper\strut
\end{minipage}\tabularnewline
\midrule
\endhead
\begin{minipage}[t]{0.11\columnwidth}\raggedright
\(P_1\)\strut
\end{minipage} & \begin{minipage}[t]{0.18\columnwidth}\raggedright
\textbf{descent}\strut
\end{minipage} & \begin{minipage}[t]{0.35\columnwidth}\raggedright
a property pushes down through a quotient\strut
\end{minipage} & \begin{minipage}[t]{0.26\columnwidth}\raggedright
every descent/blindness theorem (§4, §5)\strut
\end{minipage}\tabularnewline
\begin{minipage}[t]{0.11\columnwidth}\raggedright
\(P_2\)\strut
\end{minipage} & \begin{minipage}[t]{0.18\columnwidth}\raggedright
\textbf{representability}\strut
\end{minipage} & \begin{minipage}[t]{0.35\columnwidth}\raggedright
what a layer can admissibly express or select\strut
\end{minipage} & \begin{minipage}[t]{0.26\columnwidth}\raggedright
the SM selection layer (§4)\strut
\end{minipage}\tabularnewline
\begin{minipage}[t]{0.11\columnwidth}\raggedright
\(P_3\)\strut
\end{minipage} & \begin{minipage}[t]{0.18\columnwidth}\raggedright
\textbf{route mismatch}\strut
\end{minipage} & \begin{minipage}[t]{0.35\columnwidth}\raggedright
two construction routes to ``the same'' object fail to commute\strut
\end{minipage} & \begin{minipage}[t]{0.26\columnwidth}\raggedright
quantize-vs-curve (§5.1)\strut
\end{minipage}\tabularnewline
\begin{minipage}[t]{0.11\columnwidth}\raggedright
\(P_4\)\strut
\end{minipage} & \begin{minipage}[t]{0.18\columnwidth}\raggedright
\textbf{refinement}\strut
\end{minipage} & \begin{minipage}[t]{0.35\columnwidth}\raggedright
one layer refines another (common refinements, towers)\strut
\end{minipage} & \begin{minipage}[t]{0.26\columnwidth}\raggedright
the QM--GR parent \(L\) (§5.2), GUT\(\leftrightarrow\)SM (§4.4)\strut
\end{minipage}\tabularnewline
\begin{minipage}[t]{0.11\columnwidth}\raggedright
\(P_5\)\strut
\end{minipage} & \begin{minipage}[t]{0.18\columnwidth}\raggedright
\textbf{packaging}\strut
\end{minipage} & \begin{minipage}[t]{0.35\columnwidth}\raggedright
closure of structure into an auditable unit\strut
\end{minipage} & \begin{minipage}[t]{0.26\columnwidth}\raggedright
closure chains, record packets (§4.5, §§6.2--6.3)\strut
\end{minipage}\tabularnewline
\begin{minipage}[t]{0.11\columnwidth}\raggedright
\(P_6\)\strut
\end{minipage} & \begin{minipage}[t]{0.18\columnwidth}\raggedright
\textbf{audit}\strut
\end{minipage} & \begin{minipage}[t]{0.35\columnwidth}\raggedright
the ledger check that nothing was smuggled between layers\strut
\end{minipage} & \begin{minipage}[t]{0.26\columnwidth}\raggedright
the anti-smuggling discipline (§3)\strut
\end{minipage}\tabularnewline
\bottomrule
\end{longtable}

Two disclaimers are part of the calculus itself (and are proved, not just stated, in Foundations III): the labels form
\textbf{no algebra} --- there is no total operation \(*:\mathbb P^2 \to \mathbb P\), and a ``no total six-symbol algebra'' theorem
forbids reading interactions off the labels alone; and there is \textbf{no universal reduction claim} --- the calculus does not
assert that all of physics decomposes into six primitives. The six are role names under which finite, typed,
individually-audited judgments are filed.

On top of the primitives sits a catalog of \textbf{layer-agnostic structural laws} (``F-laws'', Foundations IV~\citep{Tsiokos_FoundationsIV}): theorem-grade
statements about quotients and closures that hold regardless of substrate, each verified formally and instantiated
across multiple sciences. The ones used in this paper (one-line forms in Appendix A): F23 (probability as a stable
measure on an unresolved fiber), F27 (conservation as orbit descent), F34 (information loss as an obstruction-set
verdict), F37 (complementarity as non-existence of a joint quotient), F39 (entropy as fiber volume), F47/F26
(fine-tuning and contingency: no value-law lands on a contingent selection), F48 (topological obstruction/gluing), F49
(common-source/nonlocal correlation normal form), F50 (vacuum/background selection: necessary vs moduli vs vacuum), F51
(unification as common refinement).

\hypertarget{what-counts-as-a-result-landing-modes}{\subsection{What counts as a result: landing modes}\label{what-counts-as-a-result-landing-modes}}

Because the calculus is built to \emph{audit} claims, every result in this paper carries one of a small set of grades --- the
\textbf{landing modes} --- and the grade is part of the result:

\begin{itemize}
\tightlist
\item
  \textbf{COMPUTE} --- a relation, number, or verdict drops out of the construction (e.g., the factorization defects of §4--5).
\item
  \textbf{SELECT} --- the machinery picks a structure out of a space, \emph{on declared conditions} (e.g., the SM gauge algebra,
  conditional on clean-separation).
\item
  \textbf{PROVE-BLIND} --- a proof that the calculus \emph{cannot} determine a quantity, which thereby certifies it as
  observed-input. This is a positive result: it demarcates. (E.g., the generation-count blindness theorem of §4.6.)
\item
  \textbf{GROUND} --- a condition not derivable within its own layer is supplied as the \emph{down-shadow of a named higher source},
  with the source and its warrant declared. A GROUND landing is conditional and is tagged as such. (E.g.,
  clean-separation grounded in record-stability, §4.5; the QM--GR common carrier grounded in semiclassical co-sourcing,
  §5.3.)
\item
  \textbf{RECOGNITION} --- a known result is shown to be a special case of one of the calculus's laws, \emph{non-circularly}: the
  law's machinery reproduces it without the result being inserted (e.g., Ryu--Takayanagi as ledger/shadow-price duality,
  §5.4).
\item
  \textbf{FORBIDDEN-RULE / DEMARCATION} --- the calculus forbids a configuration, or classifies a quantity as law-fixable
  versus contingent, with an explicit can-fail control (e.g., the predictions of §6; the Λ classification of §7.3).
\end{itemize}

The grades keep these distinctions explicit: a SELECT is not presented as a derivation, a
RECOGNITION not as a discovery, and a toy-level FORBIDDEN-RULE not as a fact about nature. Conversely ---
and this matters for Section 6 --- a forced, falsifiable prediction is not \emph{deflated} into ``recovery'' merely because its
subject matter is familiar.

\hypertarget{the-anti-smuggling-discipline}{\subsection{The anti-smuggling discipline}\label{the-anti-smuggling-discipline}}

The standing failure mode of any flexible framework is to smuggle its conclusion into its premises. The calculus's
\(P_6\) (audit) role is a set of mechanical gates against this: a quantity may not be \emph{defined} to equal the target
(anti-tautology); each hypothesis must be independently satisfiable by a structure that \emph{fails} the conclusion
(anti-circularity); imported machinery is frozen and hash-pinned so it cannot be retuned toward the answer
(frozen-machinery); and every positive claim must be accompanied by a \textbf{can-fail control} --- a configuration on which
the claim comes out false, demonstrating the test discriminates. Section 3 describes how these gates operated in
practice, including the cases where they fired against our own constructions.
 \hypertarget{method-finite-audited-carriers-adversarial-self-correction}{\section{Method: finite audited carriers, adversarial self-correction}\label{method-finite-audited-carriers-adversarial-self-correction}}

Nothing in this paper is an estimate, a fit, or a numerical simulation in the usual sense. Every result is an \textbf{exact
computation on a finite, frozen, reproducible carrier}, produced under a protocol designed to catch the framework ---
and its operators --- rigging the result. This section states the protocol; Appendix C documents the occasions it fired.

\hypertarget{finite-frozen-carriers}{\subsection{Finite frozen carriers}\label{finite-frozen-carriers}}

Each construction declares a finite carrier (a candidate-structure space of \(11{,}990\) gauge structures in Section 4;
finite tensor-network/graph geometries and mode-grammars in Section 5), and all quantities --- obstruction sets,
factorization defects, entropies, min-cuts, kernels --- are computed deterministically on it. Carriers and predicates,
once built, are \textbf{frozen}: downstream steps import them verbatim under SHA-256 pins, and a validator fails if any
imported predicate has been altered. This forbids the most common quiet failure of long programs: retuning an early
definition until a late result comes out right.

Every step ships with an independent \textbf{validator}: a script that \emph{recomputes} the headline quantities from
the frozen inputs and fails on mismatch, checks the anti-smuggling gates of §2.4, and greps the artifacts for an
overclaim blocklist (phrases such as ``derives the gauge group'', ``solves quantum gravity'', ``proves the proton is stable''
fail the build). Validator pass is necessary, never sufficient: each accepted result was additionally re-derived
independently of its own build script (Appendix B lists the re-derivations).

\hypertarget{can-fail-controls}{\subsection{Can-fail controls}\label{can-fail-controls}}

A claim with no way to come out false is not a result. Every positive claim in this paper is paired with an explicit
\textbf{control on which it fails}, computed alongside it:

\begin{itemize}
\tightlist
\item
  the SM selection's ``collapse'' verdict is paired with a broad-measure landscape control that must \emph{not} collapse
  (and an order-battery of 204 refinement orders, 0 flips);
\item
  the QM--GR directed no-go is paired with a \emph{nested} readout for which the reduction must (and does) succeed;
\item
  the holographic monogamy results are paired with a GHZ state that must (and does) violate them;
\item
  the entanglement-kernel prediction of §6.1 is paired with a boundary-anchored tree geometry on which the kernel must
  (and does) vanish;
\item
  the record-stability prediction of §6.3 is paired with single-conjunct controls proving the implication is not
  definitional (60 and 311 explicit counterexample structures).
\end{itemize}

\hypertarget{adversarial-self-correction-in-both-directions}{\subsection{Adversarial self-correction --- in both directions}\label{adversarial-self-correction-in-both-directions}}

Over the course of the program the audit gates rejected our own constructions repeatedly. Three classes of defect were
caught and killed (full post-mortems in Appendix C):

\begin{enumerate}
\def\labelenumi{\arabic{enumi}.}
\item
  \textbf{Tautologies} --- a ``coincidence'' that is an identity. \emph{Example:} an early area--entropy ``match'' computed the von
  Neumann entropy from the Schmidt spectrum on one side and the Shannon entropy of the squared singular values on the
  other --- the same number twice, machine-\(\epsilon\) residual. Rejected; the accepted version (§5.4) uses the
  \emph{geometric} min-cut, which is provably independent of the state spectrum (it is invariant under reseeding the state
  while the entropy moves).
\item
  \textbf{Circularities / hardcoding} --- a forcing whose load-bearing clause does no work. \emph{Example:} a first version of the
  monogamy-forcing result conjoined the constraint with a compatibility flag that was constantly true, and asserted
  its control outcome as literals. Rejected; the accepted version (§5.5) defines compatibility as a computed,
  monogamy-independent geometric-dual test, and the control is computed over constructed admissible sets.
\item
  \textbf{Rigging toward a desired verdict --- in either direction.} The decisive case is Prediction 1 (§6.1), which took
  four attempts: one rigged \emph{toward} the positive verdict (a bespoke graph with the invisible mode built in by hand);
  one rigged \emph{toward} the deflationary verdict (a hardcoded ``genuine dimension \(=0\)'' plus a scoring filter chosen to
  dismiss the answer); one inflated by a \textbf{degeneracy artifact} (a one-sided finite-difference Jacobian taken at a
  maximally tied uniform-weight point, where the apparent kernel is first-order visible --- residuals scaling linearly
  with step size betrayed it). All three were rejected. The accepted computation uses generic non-degenerate
  geometries, two-sided central differences, and a both-direction finite-range invisibility test --- and its numbers
  were reproduced independently of the build before acceptance.
\end{enumerate}

We call the resulting discipline \textbf{symmetric scrutiny}: a result is attacked with equal force whether it flatters the
framework, flatters the mainstream, or flatters the authors' expectations. The rejected-attempts record is given in
Appendix C. The predictions are not the first thing the machinery produced; they are what survived an attempt to break
them in both directions.

\hypertarget{adversarial-review-and-the-anti-contamination-guard}{\subsection{Adversarial review and the anti-contamination guard}\label{adversarial-review-and-the-anti-contamination-guard}}

Two construction results were stress-tested in adversarial review packets and \textbf{settled}: the ledger/shadow-price
recognition of Ryu--Takayanagi (§5.4) and the grounding of the QM--GR common-carrier premise (§5.3); the SM construction
core separately passed an adversarial review cycle (sound-with-fixes). Review packets and dispositions are in Appendix D.
These cycles are part of the audit record, not peer review or independent replication.

Finally, because the universality thesis (§8) requires that the two tracks not contaminate each other, the program ran
under an explicit \textbf{anti-contamination guard}: the two substrates were declared structurally different at the outset
(a selection/measure layer vs a co-sourcing common refinement), every SM-track step was validated against importing
co-sourcing structure (build validators fail on field-carrier tokens), and steps that pattern-matched one substrate
onto the other were rejected. Whatever machinery the two tracks share, they share because the \emph{laws} are layer-agnostic
--- not because the constructions were allowed to copy each other.

\begin{figure}[tbp]
  \centering
  \includegraphics[width=\linewidth]{figures/fig_f0_audit_pipeline.png}
  \caption{\textbf{Audit pipeline (protocol).} \emph{One-panel schematic: declare carrier \(\to\) freeze (SHA-256) \(\to\) construct \(\to\) validator recomputes \(\to\) can-fail control \(\to\) independent re-derivation \(\to\) (for reviewed construction results) adversarial review. Side channel: the rejection loop, with the three defect classes of §3.3 feeding back.}}
  \label{fig:f0}
\end{figure}
 \hypertarget{result-i-the-standard-model-as-a-selection-layer-a-demarcation}{\section{Result I --- the Standard Model as a selection layer (a demarcation)}\label{result-i-the-standard-model-as-a-selection-layer-a-demarcation}}

\hypertarget{the-problem-restated-structurally}{\subsection{The problem, restated structurally}\label{the-problem-restated-structurally}}

The Standard Model poses five ``why \textbf{this}?'' questions --- gauge group and representations, generation count and
textures, the electroweak scale, the vacuum, the UV completion. The calculus types all five as facets of \textbf{one
selection/measure layer} \(L^*\) sitting above the SM: a candidate-structure space \(W\) together with constraints and a
selection, of which the realized SM is a single point. The structural finding that frames everything else is that the
selection is \textbf{non-descending}: no function of the realized SM's own values determines the selection (computed
obstruction witnesses; a derived observable control \emph{does} descend, so the test discriminates).

This reframes the question. One does not \emph{derive} a selection from inside its own layer; one asks \textbf{which features of
the selected point are forced by structure, and which are irreducibly contingent} --- and proves the split. The honest
answer to ``why this whole SM?'' is a \emph{factorization}:

\[
\text{SM} \;=\; \underbrace{\text{[features closure forces or grounds]}}_{\text{few, listed below}}
\;\oplus\;
\underbrace{\text{[features provably invisible to closure]}}_{\text{observed-input, \emph{proved} so}} .
\]

This is the Galois genre of result: the point is the sharp, non-circular line, in both directions.

\hypertarget{unconditional-exclusion-sim-99.3-of-chiral-gauge-structures}{\subsection{\texorpdfstring{Unconditional exclusion: \(\sim 99.3\%\) of chiral gauge structures}{Unconditional exclusion: \textbackslash sim 99.3\textbackslash\% of chiral gauge structures}}\label{unconditional-exclusion-sim-99.3-of-chiral-gauge-structures}}

On a neutral carrier of \(11{,}990\) genuinely-chiral gauge structures (built token-blind: no SM labels appear in any
predicate; the realized SM-analog is report-only), the closure chain --- anomaly freedom~\citep{Bouchiat1972}, descent, corrected chirality,
mass-closure --- \textbf{excludes all but a small family}, leaving essentially the two-factor \(2|3\) class and a single-factor
\(SU(4)\) class. No condition is imposed beyond neutrality of the machinery; the exclusion is unconditional
(\textbf{COMPUTE}). Named competitors run through the same filter: \(SU(5)\) and flipped \(SU(5)\) fail at mass-closure;
Pati--Salam~\citep{PatiSalam1974}, left--right, and trinification are excluded by route-completeness within the declared component cap
(cap-conditional, stated as such); \(SO(10)\)/\(E_6\)~\citep{FritzschMinkowski1975} are outside the declared alphabet (not refuted).

\hypertarget{conditional-selection-of-mathfraksu2oplusmathfraksu3oplusmathfraku1}{\subsection{\texorpdfstring{Conditional selection of \(\mathfrak{su}(2)\oplus\mathfrak{su}(3)\oplus\mathfrak{u}(1)\)}{Conditional selection of \textbackslash mathfrak\{su\}(2)\textbackslash oplus\textbackslash mathfrak\{su\}(3)\textbackslash oplus\textbackslash mathfrak\{u\}(1)}}\label{conditional-selection-of-mathfraksu2oplusmathfraksu3oplusmathfraku1}}

One declared condition separates the survivors: \textbf{clean separation} --- the absence of confining-charged broken vectors.
In calculus terms it is the emptiness of a factorization defect: writing \(\pi_0\) for the confining readout and \(\pi_1\)
for the mass/breaking readout on the gauge bosons of a candidate structure,

\[
\text{clean separation} \iff \Delta_{\mathrm{fact}} = \varnothing ,
\]

and the witnesses of \(\Delta_{\mathrm{fact}}\neq\varnothing\) are exactly the \(X/Y\)-type colored coset bosons. Demanding
clean separation selects \(SU(2)\times SU(3)\times U(1)\) over the single-factor alternative (\textbf{SELECT}, on one named
condition). The single-factor side has since been upgraded to a constructed theorem: \emph{a single \(SU(N)\) with a stable
confining substrate has \(\Delta_{\mathrm{fact}}\neq\varnothing\) for all \(N\)} (symbolic three-step proof from the frozen
machinery; non-vacuous at \(N=4\)).

\hypertarget{the-xy-coset-spine-and-the-recovered-unification-ratios}{\subsection{\texorpdfstring{The \(X/Y\)-coset spine and the recovered unification ratios}{The X/Y-coset spine and the recovered unification ratios}}\label{the-xy-coset-spine-and-the-recovered-unification-ratios}}

The same defect object organizes a cluster of phenomena usually treated separately:

\begin{quote}
\textbf{One object, four roles.} The \(X/Y\) colored coset is simultaneously (i) the clean-separation violator, (ii) the
proton-decay mediator~\citep{Langacker1981} (its absence makes baryon number an F27 orbit-descent invariant: \(\mathcal O = 0\)), (iii) the
monopole source~\citep{Dirac1931,tHooft1974monopole} (its absence zeroes the F48 gluing obstruction), and (iv) the witness that a grand-unified parent is
a genuine refinement of the SM rather than a relabeling: \(\Delta_{\mathrm{fact}}(\mathrm{SM},\mathrm{GUT}) = 15\),
\(\Delta_{\mathrm{fact}}(\mathrm{GUT},\mathrm{SM}) = 0\).
\end{quote}

\begin{figure}[tbp]
  \centering
  \includegraphics[width=\linewidth]{figures/fig_f2_xy_coset_spine.png}
  \caption{\textbf{\(X/Y\)-coset spine.} \emph{One central node (\(\Delta_{\mathrm{fact}}\) witnesses) with four arrows --- clean-separation, proton stability (F27), no monopole (F48), GUT non-relabeling --- annotated ``absence of one object \(\Rightarrow\) all four at once.''}}
  \label{fig:f2}
\end{figure}

Constructing the F51 common-refinement parent over the SM's two routes lands the embedding numbers \textbf{as computations,
not inputs}. Over one SM generation,

\[
\sin^2\theta_W \;=\; \frac{\operatorname{Tr} T_3^2}{\operatorname{Tr} Q^2}
\;=\; \frac{2}{16/3} \;=\; \frac{3}{8},
\qquad
k_Y \;=\; \frac{5}{3} \ \ \text{(from } \operatorname{Tr} Y^2 = \tfrac{5}{6}\text{)} ,
\]

with the traces computed from the SM charge table (declared input) and the \emph{minimal-simple} parent class (declared
source). Non-circularity has a computed control: a product parent satisfying the same bare requirement yields \(3/23\),
so the value \(3/8\) genuinely depends on the declared source rather than being baked in (\textbf{GROUND}, conditional on that
source). Scope is strict: these are the \emph{embedding} ratios; the measured low-energy \(\sin^2\theta_W\) requires
renormalization-group running to a physical scale, which the toy does not contain --- and the flagship unification test
\textbf{fails} for minimal non-supersymmetric \(SU(5)\)~\citep{GeorgiGlashow1974}, consistent with experiment~\citep{SuperK2020}.

\hypertarget{grounding-the-condition-record-stability}{\subsection{Grounding the condition: record-stability}\label{grounding-the-condition-record-stability}}

Clean separation itself is not toy-derivable (the only internal grounding is circular --- caught and rejected). It is
\textbf{grounded one layer up} (\textbf{GROUND}): requiring a \emph{memory/record-stability} layer --- a stable confining substrate, the
capacity for at least two neutral records, and distinguishability --- \textbf{forces} clean separation. On the full carrier:

\[
\#\{\,\text{record-stable} \wedge \Delta_{\mathrm{fact}}\neq\varnothing\,\} \;=\; 0
\quad\text{out of } 11{,}990 .
\]

Anti-circularity is computed, not asserted: the substrate conjunct alone admits \(60\) decaying structures, the capacity
conjunct alone admits \(311\); only the conjunction lands in the clean class, and record-stability is \emph{not} extensionally
clean-separation (\(24\) record-stable structures versus \(316\) clean --- a strict directional implication). The
all-structures theorem upgrade (lemma L60\(\to\)L64) remains open and is tracked as such.

\hypertarget{the-blindness-theorems-what-is-proved-observed-input}{\subsection{\texorpdfstring{The blindness theorems: what is \emph{proved} observed-input}{The blindness theorems: what is proved observed-input}}\label{the-blindness-theorems-what-is-proved-observed-input}}

The demarcation's other half is proof-grade negatives (\textbf{PROVE-BLIND}):

\begin{itemize}
\tightlist
\item
  \textbf{Generation count.} For an anomaly-free content unit with an even per-unit \(SU(2)\)-doublet count (the SM unit
  qualifies), every gate of the frozen closure chain is invariant in \(N_{\mathrm{gen}}\): anomaly coefficients scale
  linearly (\(N\cdot u\)), the Witten parity gate~\citep{Witten1982} depends only on \((N\cdot d)\bmod 2\), and the remaining gates are
  nonempty-gated. Hence the chain passes identically for \textbf{all} \(N \ge 1\) (constructed theorem; an odd-doublet control
  \emph{is} \(N\)-sensitive, so the hypothesis is load-bearing). The calculus cannot see \(N_{\mathrm{gen}}=3\); the only handle
  is the recognition-source bound \((N-1)(N-2)/2 > 0 \iff N \ge 3\) from CP violation~\citep{KobayashiMaskawa1973} --- observed-input.
\item
  \textbf{Content.} Within the selected gauge structure, three definitionally distinct neutral shadows (integer charge
  quantization of the color-singlet spectrum; Yukawa-texture connectivity; mass-matrix rank) are all blind to the
  SM-vs-alternative content classes: the detailed matter content is a contingency-class residual (F26).
\item
  \textbf{Scales and values.} The electroweak scale is reframed (not solved) as an F47 small-selector region (\(\approx 3.96\%\)); no measured low-energy observable is derived anywhere in the track.
\end{itemize}

\hypertarget{result-sm-summary-table}{\subsection{Summary table}\label{result-sm-summary-table}}

\begin{longtable}[]{@{}lll@{}}
\toprule
\begin{minipage}[b]{0.28\columnwidth}\raggedright
question\strut
\end{minipage} & \begin{minipage}[b]{0.31\columnwidth}\raggedright
verdict (grade)\strut
\end{minipage} & \begin{minipage}[b]{0.34\columnwidth}\raggedright
basis\strut
\end{minipage}\tabularnewline
\midrule
\endhead
\begin{minipage}[t]{0.28\columnwidth}\raggedright
exclude \(\sim99.3\%\) of chiral gauge structures\strut
\end{minipage} & \begin{minipage}[t]{0.31\columnwidth}\raggedright
\textbf{COMPUTE}, unconditional\strut
\end{minipage} & \begin{minipage}[t]{0.34\columnwidth}\raggedright
closure chain on \(11{,}990\) structures\strut
\end{minipage}\tabularnewline
\begin{minipage}[t]{0.28\columnwidth}\raggedright
select \(\mathfrak{su}(2)\oplus\mathfrak{su}(3)\oplus\mathfrak{u}(1)\)\strut
\end{minipage} & \begin{minipage}[t]{0.31\columnwidth}\raggedright
\textbf{SELECT}, on clean separation\strut
\end{minipage} & \begin{minipage}[t]{0.34\columnwidth}\raggedright
\(\Delta_{\mathrm{fact}}=\varnothing\)\strut
\end{minipage}\tabularnewline
\begin{minipage}[t]{0.28\columnwidth}\raggedright
proton stability \(+\) no monopole \(+\) color separation\strut
\end{minipage} & \begin{minipage}[t]{0.31\columnwidth}\raggedright
\textbf{COMPUTE}\strut
\end{minipage} & \begin{minipage}[t]{0.34\columnwidth}\raggedright
one object: the \(X/Y\) coset, absent\strut
\end{minipage}\tabularnewline
\begin{minipage}[t]{0.28\columnwidth}\raggedright
\(\sin^2\theta_W = 3/8\), \(k_Y = 5/3\)\strut
\end{minipage} & \begin{minipage}[t]{0.31\columnwidth}\raggedright
\textbf{GROUND} (minimal-simple parent)\strut
\end{minipage} & \begin{minipage}[t]{0.34\columnwidth}\raggedright
trace ratios; product control \(3/23\)\strut
\end{minipage}\tabularnewline
\begin{minipage}[t]{0.28\columnwidth}\raggedright
clean separation itself\strut
\end{minipage} & \begin{minipage}[t]{0.31\columnwidth}\raggedright
\textbf{GROUND} (record-stability)\strut
\end{minipage} & \begin{minipage}[t]{0.34\columnwidth}\raggedright
\(0/11{,}990\) counterexamples; non-circular\strut
\end{minipage}\tabularnewline
\begin{minipage}[t]{0.28\columnwidth}\raggedright
single-factor exclusion, all \(N\)\strut
\end{minipage} & \begin{minipage}[t]{0.31\columnwidth}\raggedright
\textbf{COMPUTE} (constructed theorem)\strut
\end{minipage} & \begin{minipage}[t]{0.34\columnwidth}\raggedright
symbolic proof\strut
\end{minipage}\tabularnewline
\begin{minipage}[t]{0.28\columnwidth}\raggedright
\(N_{\mathrm{gen}}\), content, measured values\strut
\end{minipage} & \begin{minipage}[t]{0.31\columnwidth}\raggedright
\textbf{PROVE-BLIND} / observed-input\strut
\end{minipage} & \begin{minipage}[t]{0.34\columnwidth}\raggedright
all-\(N\) theorem; three blind shadows\strut
\end{minipage}\tabularnewline
\begin{minipage}[t]{0.28\columnwidth}\raggedright
measured \(\sin^2\theta_W\), couplings, masses, EW scale\strut
\end{minipage} & \begin{minipage}[t]{0.31\columnwidth}\raggedright
\textbf{not landed} (out of scope)\strut
\end{minipage} & \begin{minipage}[t]{0.34\columnwidth}\raggedright
need running/scale = experiment\strut
\end{minipage}\tabularnewline
\bottomrule
\end{longtable}

The Standard Model, on this evidence, is not one derivable object. It is a \emph{selection} with a thin forced skeleton ---
and the calculus proves both the skeleton and the thinness. Section 6.3 turns the spine of this analysis into a
falsifiable prediction.
 \hypertarget{result-ii-quantum-mechanics-general-relativity-as-a-common-refinement-a-unification}{\section{Result II --- quantum mechanics ↔ general relativity as a common refinement (a unification)}\label{result-ii-quantum-mechanics-general-relativity-as-a-common-refinement-a-unification}}

\hypertarget{locating-the-obstruction-the-route-mismatch}{\subsection{Locating the obstruction: the route mismatch}\label{locating-the-obstruction-the-route-mismatch}}

Why do QM and GR not compose? The calculus's answer is structural and exact. Call \(E_1\) the completion ``quantize'' and
\(E_2\) the completion ``curve''. Each is individually consistent --- idempotent on the carrier:

\[
E_1^2 = E_1,\qquad E_2^2 = E_2,\qquad\text{but}\qquad E_1 E_2 \neq E_2 E_1 .
\]

A theorem-grade result of the calculus (Foundations III~\citep{Tsiokos_FoundationsIII}, Thms. 9--10) shows this situation is well-posed and finitely
witnessed: the \textbf{route-mismatch defect}

\[
\Delta_3^{\mathrm{comp}}(E_1,E_2) \;=\; \{\,c : E_1E_2(c) \neq E_2E_1(c)\,\}
\]

is non-empty iff the completions fail to commute, \emph{even though each is exact alone} --- i.e., route mismatch is a real,
typed, finite defect, not a symptom of sloppiness in either route. The explicit minimal witness is elementary: on
\(C=\mathcal P(\{a,b,c\})\) with \(E_1(S)=S\cup\{b\}\) iff \(a\in S\) and \(E_2(S)=S\cup\{c\}\) iff \(b\in S\), one has
\(E_1E_2(\{a\})=\{a,b\}\) but \(E_2E_1(\{a\})=\{a,b,c\}\). On the QM--GR carrier below, the directed route mismatch computes
to \(0.5\) (normalized). This is the calculus's typing of \emph{why quantizing gravity fails}~\citep{tHooftVeltman1974,GoroffSagnotti1985}: quantize-then-curve and
curve-then-quantize are different objects, and no amount of skill in either route removes the defect of composing them.

\hypertarget{the-resolution-a-fork-not-a-ladder}{\subsection{The resolution: a fork, not a ladder}\label{the-resolution-a-fork-not-a-ladder}}

The construction models the two theories as \textbf{access quotients of one carrier}. The minimal mode-grammar has four
modes \((d_0,d_1,d_2,d_3)\) with

\[
q_{\mathrm{QM}} = (d_0,d_1,d_2), \qquad q_{\mathrm{GR}} = (d_0,d_2,d_3), \qquad L = (d_0,d_1,d_2,d_3),
\]

i.e., the theories share modes \((d_0,d_2)\), while \(d_1\) (phase/superposition data) is QM-only and \(d_3\) (curvature
readout) is GR-only. Two architectures compete:

\begin{itemize}
\tightlist
\item
  the \textbf{ladder} (the emergent-spacetime direction~\citep{VanRaamsdonk2010}): GR is a lawful quotient of QM, \(q_{\mathrm{GR}} = \phi \circ q_{\mathrm{QM}}\) for some \(\phi\);
\item
  the \textbf{fork}: both are quotients of the common refinement \(L\).
\end{itemize}

The computation is decisive on the carrier. The directed ladder \textbf{fails}: \(\Delta_{\mathrm{fact}}(q_{\mathrm{QM}}, q_{\mathrm{GR}})\) has \(8\) witness pairs (row residual \(0.577\), route mismatch \(0.5\)) --- and the \emph{strongest} ladder, which
adjoins to QM an RT-shadow functional \(A_{\mathrm{RT}} = d_0 + 2d_2\) of its accessible modes, fails identically (defect
\(8\)): a function of QM-visible modes cannot recover \(d_3\). The \textbf{fork closes}: \(L \to q_{\mathrm{QM}}\) and \(L \to q_{\mathrm{GR}}\) both factor exactly (residuals \(0\)). The can-fail control flips: replacing GR by a readout genuinely
nested in QM, \(q_{\mathrm{GR}}^{\mathrm{nested}} = (d_0,d_2,d_1)\), the ladder defect drops to \(0\) --- the no-go tracks
real structure, not the test apparatus. Supporting theorems sharpen this into a program-level statement: a directed
reduction \emph{and} a fused ``quantum-gravity object'' both fail in the grammar (\(T_{\mathrm{QG\text{-}NoGo}}\)), while \(L\) is
the \textbf{unique minimal admissible common refinement} (\(T_{\mathrm{QGR\text{-}Unique}}\), type- and instance-uniqueness by
adversarial defeat, basis-independent, stable across carriers).

Scope, stated plainly: the ladder no-go covers the \textbf{weak-RT} direction (area-type functionals of accessible
entanglement data); \emph{strong} bulk reconstruction --- \(d_3\) recoverable from richer quantum data --- is not refuted here and
is exactly what Section 6.1 turns into a falsifiable test.

\begin{figure}[tbp]
  \centering
  \includegraphics[width=\linewidth]{figures/fig_f3_fork_rt.png}
  \caption{\textbf{Fork resolution and RT ledger recognition.} \emph{Left: the failed ladder (QM \(\to\) GR arrow struck through, defect 8) vs the closed fork (\(L\) above, two clean arrows down, residuals 0; nested control shown flipping). Right: panel for §5.4 --- a boundary-anchored flow network with the minimal cut highlighted; caption ``area \(=\) shadow price of the entanglement flow (max-flow/min-cut \(=\) LP duality).''}}
  \label{fig:f3}
\end{figure}

\hypertarget{grounding-the-premise}{\subsection{Grounding the premise}\label{grounding-the-premise}}

The fork is conditional on the \textbf{common-carrier premise} --- that QM and GR may be modeled as access quotients of one
carrier at all. The calculus does not let that premise pass silently. A door-test against the framework's own laws
(F37, F51, FoEC, SAU) shows each supplies the \emph{form} of a shared carrier but none \emph{forces} its existence
(form-not-existence in every case). The premise is therefore landed in \textbf{GROUND} mode: supplied as the down-shadow of
a named physical source --- \textbf{semiclassical co-sourcing}: one field \(\psi\) carries both the Born readout
\(\rho = |\psi|^2\) and the stress-energy \(T[\psi]\) that sources geometry, extended through a computed self-consistent
back-reaction \(V[\psi] = V_0 + \kappa\, T_{00}[\psi]\) (fixed-point residual \(3.5\times 10^{-16}\) at \(\kappa=0.3\); a
runaway control at \(\kappa=30\) fails, so the consistency is not automatic). Two can-fail controls give the landing
teeth: a complementary (non-commuting) access pair admits \textbf{no} joint quotient (commutator residual \(0.707\)), so the
laws do not manufacture carriers; and removing co-sourcing removes the warrant (stress residual \(0.381\)). This grounding
was stress-tested in an adversarial review cycle and \textbf{settled} as an honest recognition-source landing (Appendix D). The
named residual is real: the warrant reaches matter-sourced geometry, not the free gravitational sector or the deep
quantum-gravity regime.

\hypertarget{recognition-ryutakayanagi-as-ledgershadow-price-duality}{\subsection{Recognition: Ryu--Takayanagi as ledger/shadow-price duality}\label{recognition-ryutakayanagi-as-ledgershadow-price-duality}}

Within the fork, the calculus \emph{recognizes} the central result of holographic entanglement --- non-circularly. The claim
(a \textbf{RECOGNITION} landing, adversarially reviewed and settled): the Ryu--Takayanagi relation is a special case of the
calculus's currency--constraint law --- \emph{a conserved ledger quantity in one layer reappears as the shadow price of its
constraint in the layer above}. Concretely, on tensor-network carriers the entanglement entropy obeys the saturable
bound

\[
S(A) \;\le\; \mathrm{mincut}(A),
\]

and max-flow/min-cut~\citep{FordFulkerson1956,FreedmanHeadrick2017} \textbf{is} linear-programming duality: the minimal cut (the ``area'') is the optimal dual variable --- the
shadow price --- of the boundary entanglement flow. The construction earns the recognition rather than assuming it:
\(S(A)\) is computed from contracted random-tensor states (not defined as a cut), saturation \emph{emerges} with bond
dimension (\(S/\mathrm{mincut} = 0.729 \to 0.904 \to 0.941\) for \(D=2,3,4\), never exact at finite \(D\) --- exactness would
be the tautology signature and fails the validator), strict-gap controls exist, and the geometric side is
seed-invariant while the entropic side moves (\(\mathrm{area} = \log 3\) constant across states; \(S_{\mathrm{Born}} = 1.0970, 1.0956, 1.0980\)). Two independently checkable consequences computed on the same carrier: the holographic
entropy cone's monogamy facet~\citep{HaydenHeadrickMaloney2013},

\[
I_3(A{:}B{:}C) = S_A + S_B + S_C + S_{ABC} - S_{AB} - S_{AC} - S_{BC} \;\le\; 0
\]

for min-cut-realizable states, with the GHZ control~\citep{GHZ1989} violating it (\(I_3 = +\log 2\)); and a discrete linearized-Einstein~\citep{Faulkner2014}
\textbf{consistency condition} --- demanding RT-consistency of a perturbed geometry across all \(38\) regions of a \(14\)-edge
carrier overdetermines the response (rank \(10\), cokernel \(28\)): a geometric perturbation satisfies it (residual
\(9\times10^{-17}\)), a generic state perturbation violates it (\(0.042\)). Two unification-grade refinements complete the
picture: the quantum (Born/F23) and gravitational (area/F39) ledgers carry \textbf{one composition law} (the monogamy
combiner --- shared on co-sourced carriers, broken by GHZ), and the F51 parent \emph{forces} that combiner on admissible
co-readouts while QM alone does not (a computed geometric-dual compatibility test; the forcing disappears under a
broken-compatibility control).

We emphasize the grade: this \textbf{recovers and organizes} known holographic structure under one law --- it does not derive
holography~\citep{Maldacena1998}, the \(1/4G\) coefficient~\citep{Bekenstein1973}, or the continuum Einstein equation (the discrete condition above is the extent
of the claim; the continuum limit is an open, precisely-typed upgrade).

\hypertarget{result-qmgr-summary-table}{\subsection{Summary table}\label{result-qmgr-summary-table}}

\begin{longtable}[]{@{}lll@{}}
\toprule
\begin{minipage}[b]{0.29\columnwidth}\raggedright
question\strut
\end{minipage} & \begin{minipage}[b]{0.32\columnwidth}\raggedright
verdict (grade)\strut
\end{minipage} & \begin{minipage}[b]{0.32\columnwidth}\raggedright
basis\strut
\end{minipage}\tabularnewline
\midrule
\endhead
\begin{minipage}[t]{0.29\columnwidth}\raggedright
why QM and GR don't compose\strut
\end{minipage} & \begin{minipage}[t]{0.32\columnwidth}\raggedright
\textbf{COMPUTE} --- route mismatch \(\Delta_3^{\mathrm{comp}}\neq\varnothing\)\strut
\end{minipage} & \begin{minipage}[t]{0.32\columnwidth}\raggedright
non-commuting completions; directed mismatch \(0.5\)\strut
\end{minipage}\tabularnewline
\begin{minipage}[t]{0.29\columnwidth}\raggedright
GR a quotient of QM (emergent spacetime, weak form)?\strut
\end{minipage} & \begin{minipage}[t]{0.32\columnwidth}\raggedright
\textbf{COMPUTE} --- no (defect \(8\); control flips)\strut
\end{minipage} & \begin{minipage}[t]{0.32\columnwidth}\raggedright
ladder vs fork\strut
\end{minipage}\tabularnewline
\begin{minipage}[t]{0.29\columnwidth}\raggedright
QM, GR siblings under one parent?\strut
\end{minipage} & \begin{minipage}[t]{0.32\columnwidth}\raggedright
\textbf{COMPUTE}, conditional --- fork closes (\(0/0\)); parent unique\strut
\end{minipage} & \begin{minipage}[t]{0.32\columnwidth}\raggedright
\(T_{\mathrm{QGR\text{-}Unique}}\)\strut
\end{minipage}\tabularnewline
\begin{minipage}[t]{0.29\columnwidth}\raggedright
the common-carrier premise\strut
\end{minipage} & \begin{minipage}[t]{0.32\columnwidth}\raggedright
\textbf{GROUND} --- semiclassical co-sourcing (review-settled)\strut
\end{minipage} & \begin{minipage}[t]{0.32\columnwidth}\raggedright
door-test + two controls\strut
\end{minipage}\tabularnewline
\begin{minipage}[t]{0.29\columnwidth}\raggedright
Ryu--Takayanagi\strut
\end{minipage} & \begin{minipage}[t]{0.32\columnwidth}\raggedright
\textbf{RECOGNITION} --- ledger/shadow-price duality (review-settled)\strut
\end{minipage} & \begin{minipage}[t]{0.32\columnwidth}\raggedright
LP duality; emergent saturation; MMI; discrete Einstein\strut
\end{minipage}\tabularnewline
\begin{minipage}[t]{0.29\columnwidth}\raggedright
quantum gravity solved / continuum Einstein derived\strut
\end{minipage} & \begin{minipage}[t]{0.32\columnwidth}\raggedright
\textbf{not landed}\strut
\end{minipage} & \begin{minipage}[t]{0.32\columnwidth}\raggedright
named open upgrades\strut
\end{minipage}\tabularnewline
\bottomrule
\end{longtable}

Section 6.1 pushes the fork to its sharpest consequence --- and stakes the falsification.
 \hypertarget{the-three-falsifiable-predictions}{\section{The three falsifiable predictions}\label{the-three-falsifiable-predictions}}

The results of Sections 4--5 recover, demarcate, and unify. This section presents the three cases where the calculus
\textbf{contradicts a prevailing program and attaches a falsification condition to the result}. All three
predictions are \emph{forced} --- each is the sharp form of its track's central structural result, not an add-on --- and each
was computed under the symmetric-scrutiny protocol of Section 3 (surviving rejected riggings toward \emph{and} against it;
Appendix C). Predictions 1--2 are the sharp forms of the QM--GR fork; Prediction 3, of the SM coset spine. One
record/memory layer ties the prediction set at the grammar level, while the physical claims remain track-specific.

\begin{center}\rule{0.5\linewidth}{0.5pt}\end{center}

\hypertarget{prediction-1-entanglement-does-not-determine-geometry-strong-it-from-qubit-fails}{\subsection{Prediction 1 --- Entanglement does not determine geometry (strong it-from-qubit fails)}\label{prediction-1-entanglement-does-not-determine-geometry-strong-it-from-qubit-fails}}

\textbf{The contested program.} A prominent line in quantum gravity (``it-from-qubit'', ER=EPR, entanglement-builds-geometry)~\citep{VanRaamsdonk2010,MaldacenaSusskind2013}
holds that bulk spacetime is \emph{fully determined by} boundary entanglement: given the complete entanglement data of the
boundary state, the bulk geometry is reconstructible. The fork of §5.2 forces the opposite: the gravitational readout
\(d_3\) is a sibling of the entanglement modes, not a function of them, so \emph{some} geometric data must be invisible to
\emph{all} entanglement measures.

\textbf{The sharp form.} Let \(w \in \mathbb R^{n}_{>0}\) be the bulk geometry (edge weights/capacities of a holographic
graph carrier) and let the \textbf{entanglement fingerprint} \(\mathcal E(w)\) be the complete boundary data: all
\(2^8 - 2 = 254\) boundary-region min-cut entropies, all \(3{,}025\) mutual informations, all \(7{,}770\) tripartite
informations \(I_3\) --- \(11{,}049\) features. The prediction concerns the kernel of the reconstruction map:

\[
J \;=\; \frac{\partial\, \mathcal E}{\partial\, w}
\qquad\text{and}\qquad
\boxed{\;\dim \ker J \;>\; 0\;}
\]

robustly on non-degenerate geometries --- together with the existence of \textbf{finite, two-sided invisible deformations}:
single bulk capacities \(w_e\) that can be moved by finite \(\pm\delta\) with \emph{every one} of the \(11{,}049\) entanglement
features unchanged, while named bulk-geometric quantities change.

\textbf{The computation (independently reproduced).} On the frozen 14-edge holographic carrier of §5.4, at generic
(tie-broken) geometries, with two-sided central differences, across three seeds:

\begin{longtable}[]{@{}llll@{}}
\toprule
\begin{minipage}[b]{0.09\columnwidth}\raggedright
seed\strut
\end{minipage} & \begin{minipage}[b]{0.14\columnwidth}\raggedright
\(\operatorname{rank} J\)\strut
\end{minipage} & \begin{minipage}[b]{0.15\columnwidth}\raggedright
\(\dim\ker J\)\strut
\end{minipage} & \begin{minipage}[b]{0.52\columnwidth}\raggedright
clean shadow edges (both-direction, finite range)\strut
\end{minipage}\tabularnewline
\midrule
\endhead
\begin{minipage}[t]{0.09\columnwidth}\raggedright
101\strut
\end{minipage} & \begin{minipage}[t]{0.14\columnwidth}\raggedright
9\strut
\end{minipage} & \begin{minipage}[t]{0.15\columnwidth}\raggedright
\textbf{5}\strut
\end{minipage} & \begin{minipage}[t]{0.52\columnwidth}\raggedright
\textbf{2}\strut
\end{minipage}\tabularnewline
\begin{minipage}[t]{0.09\columnwidth}\raggedright
202\strut
\end{minipage} & \begin{minipage}[t]{0.14\columnwidth}\raggedright
9\strut
\end{minipage} & \begin{minipage}[t]{0.15\columnwidth}\raggedright
\textbf{5}\strut
\end{minipage} & \begin{minipage}[t]{0.52\columnwidth}\raggedright
\textbf{2}\strut
\end{minipage}\tabularnewline
\begin{minipage}[t]{0.09\columnwidth}\raggedright
303\strut
\end{minipage} & \begin{minipage}[t]{0.14\columnwidth}\raggedright
9\strut
\end{minipage} & \begin{minipage}[t]{0.15\columnwidth}\raggedright
\textbf{5}\strut
\end{minipage} & \begin{minipage}[t]{0.52\columnwidth}\raggedright
\textbf{2}\strut
\end{minipage}\tabularnewline
\bottomrule
\end{longtable}

An explicit witness pair: two geometries differing by \(\delta = \pm 0.15\) on a shadow edge have \emph{identical} fingerprints
(residual \(< 10^{-9}\) across all \(11{,}049\) features, both signs) while the bulk interior volume differs by \(0.15\). The
\textbf{can-fail control} is essential and passes: on a boundary-anchored tree (every edge on some boundary min-cut) the
kernel is exactly zero --- entanglement \emph{does} determine that geometry --- so the test can return the it-from-qubit verdict
and, on the holographic carrier, does not. A degenerate-geometry artifact that initially inflated the kernel (one-sided
differences at a maximally tied point) was caught and removed; the reported numbers are the artifact-free ones
(Appendix C.3).

\textbf{The prediction, stated for the community.}

\begin{quote}
On any non-trivial holographic / tensor-network geometry, the map from the complete boundary-entanglement fingerprint
(all subregion entropies, mutual informations, and multipartite informations) to the bulk geometry has a \textbf{non-trivial
kernel}: there is a robust sector of bulk geometric degrees of freedom invisible to every entanglement measure.
Complete entanglement-only bulk reconstruction is therefore impossible; any reconstruction program must add
non-entanglement boundary data.

\textbf{Falsifier:} exhibit a non-trivial (non-tree, multi-path) holographic or tensor-network geometry whose
entanglement\(\to\)geometry map is injective --- central-difference kernel zero and no two-sided finite invisible
deformation. That would refute the fork's sharp form.
\end{quote}

\textbf{Grade and overlap.} FORBIDDEN-RULE at toy level; conditional on the grounded premise of §5.3; not a claim
proven about nature. The \emph{phenomenon} overlaps the known ``entanglement shadow'' of holography~\citep{Freivogel2015} --- the calculus's
contribution is to \emph{force} it from a structural law (the fork), compute its dimension, exhibit a two-sided witness, and
attach a falsification condition; the overlap is a fact, not a retraction.

\begin{center}\rule{0.5\linewidth}{0.5pt}\end{center}

\hypertarget{prediction-2-gravity-does-not-mediate-entanglement-bmv-null}{\subsection{Prediction 2 --- Gravity does not mediate entanglement (BMV-null)}\label{prediction-2-gravity-does-not-mediate-entanglement-bmv-null}}

\textbf{The contested program.} The most anticipated near-term probe of quantum gravity is the BMV proposal (Bose et al.;
Marletto--Vedral)~\citep{Bose2017,MarlettoVedral2017}: place two masses in spatial superpositions, let only gravity couple them, and test whether they
become \emph{entangled}. The widely-stated inference is that gravitationally-induced entanglement would witness a \emph{quantum}
(coherent) gravitational mediator --- a graviton-like degree of freedom. The fork of §5.2 forces the opposite: the
gravitational arm is a sibling readout that does not have access to the quantum phase mode and cannot carry it
coherently between parties.

\textbf{The sharp form.} Two matter parties \(A,B\) are prepared as independent path superpositions \(|+\rangle_A|+\rangle_B\)
(initial entanglement zero, computed: \(6.2\times10^{-18}\)). A mediator couples to them \emph{only through their
geometry-visible modes}. The fork grammar fixes the mediator's character --- and this is the load-bearing derivation, not
a modeling choice: \(q_{\mathrm{GR}} = (d_0,d_2,d_3)\) \textbf{excludes} the QM-only phase mode \(d_1\), and \(T_{\mathrm{QG\text{-} NoGo}}\) forecloses the fused amplitude-over-geometry object --- so the admissible mediator is a \textbf{classically-indexed
record channel} (one mediator state per geometry-visible configuration; a record, not a superposition). The prediction
is that such a channel generates no entanglement:

\[
\boxed{\ \mathcal N_{AB}^{\mathrm{fork}} = 0,\qquad \mathrm{CHSH}^{\mathrm{fork}} \le 2,\qquad
\text{local coherence fully damped}\ }
\]

with the \emph{forbidden} coherent mediator built only as a can-fail control that \textbf{must} entangle.

\textbf{The computation (verified against closed-form QM).} On two path qubits with a mediator of record-dimension 4:

\begin{longtable}[]{@{}lrrr@{}}
\toprule
\begin{minipage}[b]{0.18\columnwidth}\raggedright
channel\strut
\end{minipage} & \begin{minipage}[b]{0.24\columnwidth}\raggedleft
negativity \(\mathcal N\)\strut
\end{minipage} & \begin{minipage}[b]{0.24\columnwidth}\raggedleft
CHSH\strut
\end{minipage} & \begin{minipage}[b]{0.24\columnwidth}\raggedleft
local-coherence damping\strut
\end{minipage}\tabularnewline
\midrule
\endhead
\begin{minipage}[t]{0.18\columnwidth}\raggedright
fork-admissible \textbf{record} channel\strut
\end{minipage} & \begin{minipage}[t]{0.24\columnwidth}\raggedleft
\(\mathbf{0}\)\strut
\end{minipage} & \begin{minipage}[t]{0.24\columnwidth}\raggedleft
\(\mathbf{0}\)\strut
\end{minipage} & \begin{minipage}[t]{0.24\columnwidth}\raggedleft
\(1.0\) (full decoherence)\strut
\end{minipage}\tabularnewline
\begin{minipage}[t]{0.18\columnwidth}\raggedright
forbidden \textbf{coherent} control (CPhase \(\pi/2\))\strut
\end{minipage} & \begin{minipage}[t]{0.24\columnwidth}\raggedleft
\(\tfrac{1}{2\sqrt2} = 0.3536\)\strut
\end{minipage} & \begin{minipage}[t]{0.24\columnwidth}\raggedleft
\(\sqrt 6 = 2.4495\)\strut
\end{minipage} & \begin{minipage}[t]{0.24\columnwidth}\raggedleft
\(1-\tfrac{1}{\sqrt2}=0.2929\)\strut
\end{minipage}\tabularnewline
\bottomrule
\end{longtable}

with the local bound enumerated to \(2\) (16 deterministic strategies). The fork channel is exact dephasing in the joint
path basis (output \(\mathbb 1/4\)): decoherence \emph{without} entanglement. The control entangles past the classical bound,
so the test discriminates --- a channel of this finite model \emph{can} produce gravitational entanglement; the fork's channel
does not. All four numbers match closed form (\(\mathcal N = \mathcal C/2\), \(\mathrm{CHSH}=2\sqrt{1+\mathcal C^2}\) for
concurrence \(\mathcal C\)). The \textbf{mean-field trap is avoided by construction}: the channel restriction is derived from
the frozen access structure, \emph{not} from the semiclassical back-reaction \(V=V_0+\kappa T_{00}\) (an expectation-value
channel through which ``no entanglement'' would be a LOCC tautology); the validator fails on any \(T_{00}\)/expectation
coupling.

\textbf{The prediction, stated for the community.}

\begin{quote}
Under the fork/co-sourcing reading, the gravitational interaction is a classically-indexed record channel:
gravitationally-coupled masses in spatial superposition \textbf{decohere without becoming entangled}. BMV-type experiments
will return a \textbf{null} entanglement result, accompanied by which-path decoherence.

\textbf{Falsifier:} a confirmed BMV-type detection of gravitationally-induced entanglement between masses. (Secondary,
practically out of reach: a confirmed quantization of the free gravitational field / single-graviton detection --- by
Dyson's argument~\citep{Dyson2013} this is not a realistic test, so BMV is the operative one.) Either of the \emph{first} kind refutes the
fork reading of gravity directly.
\end{quote}

\textbf{Grade and overlap.} FORBIDDEN-RULE at toy level; conditional on the grounded premise of §5.3; not a claim
proven about nature. The lemma that an expectation-value/classical channel cannot create entanglement is standard
(LOCC; the Kafri--Taylor--Milburn classical-channel-gravity model~\citep{Kafri2014} is in this camp) --- the calculus's contribution is to
\emph{force} gravity onto that side of the dichotomy from the frozen fork and no-go (with the foreclosed coherent control
demonstrating it was not assumed), and to tie the gravitational channel's classicality to the \emph{same record layer} that
grounds §4 and §6.3. This is the most near-term falsifiable of the three predictions.

\hypertarget{prediction-3-record-stability-forbids-proton-decay-and-magnetic-monopoles-contra-grand-unification}{\subsection{Prediction 3 --- Record-stability forbids proton decay and magnetic monopoles (contra grand unification)}\label{prediction-3-record-stability-forbids-proton-decay-and-magnetic-monopoles-contra-grand-unification}}

\textbf{The contested program.} Grand unification generically predicts that the proton decays~\citep{Langacker1981} (lifetimes \(\sim 10^{34\text{–}36}\) yr in surviving models) and that monopoles exist: both are mediated by the \(X/Y\)-type coset bosons
of the unified parent. The spine of §4 forces the opposite --- not ``decay is slow'' but ``decay is structurally forbidden,
for a reason no other framework names.''

\textbf{The sharp form.} Define, on the frozen \(11{,}990\)-structure carrier of §4: \(\mathrm{RS}(C)\) --- structure \(C\) is
\textbf{record-stable} (stable confining substrate \(\wedge\) capacity for \(\ge 2\) neutral records \(\wedge\)
distinguishability); and the branch split by the factorization defect, \(\Delta_{\mathrm{fact}}(C) = \varnothing\)
(clean: baryon number descends, F27 obstruction \(0\); monopole gluing obstruction, F48, \(=0\)) versus
\(\Delta_{\mathrm{fact}}(C) \neq \varnothing\) (breaking: \(X/Y\) coset present; proton decay and monopoles mediated if
realized). The prediction is the implication

\[
\boxed{\;\mathrm{RS}(C) \;\Longrightarrow\; \Delta_{\mathrm{fact}}(C) = \varnothing\;}
\qquad\text{i.e.}\qquad
\#\{\, C : \mathrm{RS}(C) \wedge \Delta_{\mathrm{fact}}(C)\neq\varnothing \,\} = 0 .
\]

\textbf{The computation (independently reproduced).}

\begin{longtable}[]{@{}lrr@{}}
\toprule
\begin{minipage}[b]{0.25\columnwidth}\raggedright
\strut
\end{minipage} & \begin{minipage}[b]{0.33\columnwidth}\raggedleft
clean (\(\Delta_{\mathrm{fact}}=\varnothing\))\strut
\end{minipage} & \begin{minipage}[b]{0.33\columnwidth}\raggedleft
breaking (\(\Delta_{\mathrm{fact}}\neq\varnothing\))\strut
\end{minipage}\tabularnewline
\midrule
\endhead
\begin{minipage}[t]{0.25\columnwidth}\raggedright
record-stable\strut
\end{minipage} & \begin{minipage}[t]{0.33\columnwidth}\raggedleft
24\strut
\end{minipage} & \begin{minipage}[t]{0.33\columnwidth}\raggedleft
\textbf{0}\strut
\end{minipage}\tabularnewline
\begin{minipage}[t]{0.25\columnwidth}\raggedright
not record-stable\strut
\end{minipage} & \begin{minipage}[t]{0.33\columnwidth}\raggedleft
292\strut
\end{minipage} & \begin{minipage}[t]{0.33\columnwidth}\raggedleft
11,674\strut
\end{minipage}\tabularnewline
\bottomrule
\end{longtable}

Non-circularity is computed, in two ways. First, the hypothesis is not the conclusion in disguise:
record-stability and clean-separation are \emph{different} sets (\(24 \subsetneq 316\) --- most clean structures cannot bear
records; the implication is strictly directional). Second, each conjunct of record-stability \emph{alone} admits decaying
structures --- \(60\) substrate-only and \(311\) capacity-only breaking-branch witnesses --- so only the conjunction forces the
clean branch. The implication has content; it is not definitional.

\textbf{The prediction, stated for the community.}

\begin{quote}
The capacity of a universe to bear stable records is structurally tied to baryon-number conservation and the absence
of magnetic monopoles: any record-bearing structure lies in the clean branch, in which the coset that would mediate
proton decay and monopoles is absent. The proton of a record-bearing universe is \textbf{exactly stable} --- not long-lived
--- and no monopole exists. The distinctive content is the \textbf{link}: no other framework ties record-bearing capacity
to baryon conservation; grand unification predicts the opposite on both counts.

\textbf{Falsifier:} a confirmed proton-decay event (e.g., Hyper-Kamiokande~\citep{HyperK2025}) or a confirmed magnetic monopole. Our universe
manifestly bears stable records; either observation severs the predicted link.
\end{quote}

\textbf{Grade and scope.} FORBIDDEN-RULE / DEMARCATION at enumeration strength: \(0/11{,}990\) on the frozen
carrier, with the all-structures theorem upgrade (lemma L60\(\to\)L64) open; conditional on record-stability as the
grounding source (§4.5); frame-transfer-limited --- this does not \emph{prove} the physical proton stable. The surface
statement ``the proton does not decay'' coincides with the unadorned Standard Model; the falsifiable novelty is the
\emph{structural reason} and the \emph{link}, which GUTs contradict.

\begin{center}\rule{0.5\linewidth}{0.5pt}\end{center}

\hypertarget{what-makes-these-predictions-rather-than-retrodictions}{\subsection{What makes these predictions rather than retrodictions}\label{what-makes-these-predictions-rather-than-retrodictions}}

Three properties, shared by all three. \textbf{(i) Forced:} each is the sharp form of its track's central theorem (the fork,
for Predictions 1--2; the coset spine + grounding, for Prediction 3), not a tunable add-on --- the machinery that produces
it was frozen before the question was asked. \textbf{(ii) Contrarian:} each contradicts a live mainstream program (strong
it-from-qubit; the BMV-positive expectation; grand unification), so agreement with existing data cannot have been fitted
in. \textbf{(iii) Falsifiable at stated cost:} each names a concrete observation or construction that kills it --- and in every
case the falsifier would also destroy the track's central result.
Prediction 2 is the nearest-term: BMV-class experiments are funded and advancing, and a result is plausible within the
decade.

\begin{figure}[tbp]
  \centering
  \includegraphics[width=\linewidth]{figures/fig_f4_predictions.png}
  \caption{\textbf{Three prediction panels.} \emph{Left (Prediction 1): two bulk graphs with visibly different interior weights, identical boundary-entanglement fingerprint table beside each; a struck-through arrow ``entanglement \(\Rightarrow\) geometry.'' Center (Prediction 2): two superposed masses, a gravitational channel drawn as a meter/record (not a wavy coherent line), output state \(\mathbb 1/4\); ``decoherence, \(\mathcal N = 0\)''; a struck-through ``entanglement''. Right (Prediction 3): the \(2\times2\) record-stable \(\times\) branch table as a population map --- the record-stable column entirely inside the clean branch; arrows to ``no proton decay (F27 \(=0\))'' and ``no monopole (F48 \(=0\))''; a struck-through GUT arrow labeled ``\(X/Y\) mediation.'' A thin band beneath all three: ``one record/memory layer'' (the cross-track tie, §8).}}
  \label{fig:f4}
\end{figure}
 \hypertarget{breadth-one-calculus-three-more-deep-problems}{\section{Breadth: one calculus, three more deep problems}\label{breadth-one-calculus-three-more-deep-problems}}

The universality thesis (§8) gains most of its force from breadth: the \emph{same} obstruction/defect/moduli machinery that
produced Sections 4--6, applied without modification, classifies three further foundational puzzles. Each result below
is a computed forbidden-rule or demarcation with its own can-fail control. None is new physics --- each recovers or
classifies known structure --- and that is their role: they show the grammar is one object, not a per-problem retrofit.

\hypertarget{bell-nonlocality-the-common-carrier-is-not-a-hidden-variable}{\subsection{Bell nonlocality: the common carrier is not a hidden variable}\label{bell-nonlocality-the-common-carrier-is-not-a-hidden-variable}}

The common-carrier premise of §5.3 invites an immediate objection: a single carrier \(\psi\) underlying both readouts
\emph{sounds like} a local hidden variable --- and local hidden variables are experimentally dead. The calculus's F49 normal
form turns the objection into a computation. For a Bell pair with standard CHSH settings~\citep{Bell1964,CHSH1969} (\(a_0 = Z\), \(a_1 = X\),
\(b_{0,1} = (Z \pm X)/\sqrt2\)), the correlation table is computed from the state, giving

\[
\mathrm{CHSH}_{\mathrm{quantum}} = 2\sqrt2 \approx 2.8284 ,
\]

while the local bound is \textbf{enumerated} --- all \(16\) deterministic local strategies, \(\max = 2\) exactly. The F49
common-source test (does any setting-independent source quotient factor both readouts?) therefore fails --- the
correlation is a \emph{statused nonlocal residual}, Bell's theorem recovered as a normal form. The can-fail control
discriminates: a classical mixture (\(\mathrm{CHSH} = \sqrt2\)) \textbf{is} locally explainable, with an explicit 16-row
hidden-variable table reproducing its correlations to \(1.8\times10^{-15}\).

The calculus-specific content is the reconciliation: F49's ``common source'' must be an \emph{admissible access quotient}, and F37
(complementarity as non-joint-access) forecloses exactly that for the non-commuting settings --- the computed setting
commutators (\(0.707\)) coincide with the carrier's complementary-pair control. The carrier lives \textbf{above} the access
structure; it is not, and cannot be, a local hidden variable. \textbf{Forbidden rule:} \emph{any reading of a common-carrier
framework in which the carrier functions as a local hidden variable is excluded} --- the premise survives Bell precisely
because the joint quotient it would require does not exist.

\hypertarget{black-hole-information-loss-is-a-property-of-the-readout-not-the-carrier}{\subsection{Black-hole information: loss is a property of the readout, not the carrier}\label{black-hole-information-loss-is-a-property-of-the-readout-not-the-carrier}}

The F34 information-loss normal form asks of any transport \(\tau: H_0 \to H_1\) with recovery readout \(\rho\) and source
information \(\sigma\): is the obstruction set

\[
\mathcal O_\rho \;=\; \{(h,h') : \rho(\tau h) = \rho(\tau h') \ \wedge\ \sigma(h) \neq \sigma(h')\}
\]

empty (recoverable) or not (loss)? Instantiated on the holographic carrier with \(\tau\) the bulk\(\to\)boundary
contraction, the verdict is exact and two-sided. The preimage classes of the full boundary state are \textbf{precisely the
internal gauge orbits}: \(\operatorname{rank} = 27\) = internal dimension (computed), and for any two preimages an
explicit gauge transformation \(G\) connecting them is solved for and verified (residual \(3\times10^{-15}\)). Hence at the
\textbf{full boundary}, the only lost interior data is gauge --- every gauge-invariant (physical) interior observable is
recoverable: loss is \emph{F34-legitimate}. At a \textbf{partial readout} (the reduced state on half the boundary --- the
``early radiation only'' analog), loss is real: an exact witness applies a unitary on the unobserved half, leaving the
observed reduced state identical (\(5\times10^{-16}\)) while gauge-invariant physical observables shift (\(\approx 10^{-2}\)). \textbf{Forbidden rule:} \emph{the holographic carrier forbids physical-information loss at the full-boundary
recovery; loss is forced at partial/coarse readouts} --- in this toy, the information paradox~\citep{Hawking1976} is a property of the
readout, not of the carrier. (No dynamics, no evaporation, no Page curve~\citep{Page1993} is claimed.)

\hypertarget{the-cosmological-constant-no-in-layer-value-law}{\subsection{The cosmological constant: no in-layer value-law}\label{the-cosmological-constant-no-in-layer-value-law}}

The F50 vacuum/background-selection law classifies a background quantity as \textbf{necessary} (closure fixes one value),
\textbf{moduli} (closure is blind --- contingent), or \textbf{vacuum-selected} (a singleton selector). The carrier of §5.3 contains
the cosmological-constant analog~\citep{Weinberg1989} explicitly: the matter-independent background term in \(V[\psi] = V_0 + \kappa\, T_{00}[\psi]\). Sweeping it through the frozen closure chain: \(15\) of \(18\) background values pass (degeneracy \(15\), no
singleton selector; the pass set is bounded and non-contiguous), so the background is \textbf{F50 bounded-moduli} --- the
closure does not select its value. The discrimination control is the coupling \(\kappa\), which the same chain \emph{does}
constrain (a computed admissibility boundary: \(\kappa \le 10\) converges, \(\kappa \ge 15\) runs away) --- so the audit
distinguishes closure-constrained parameters from genuine moduli; it is not an everything-is-contingent verdict.
\textbf{Demarcation:} \emph{no in-layer \(\Lambda\) value-law exists; any claimed in-layer derivation of \(\Lambda\)'s value must
smuggle a selection; the in-layer content is at most an admissibility window.} This is the same F47/F26 demarcation
pattern the SM track exhibits for its contingent quantities --- the cross-track recurrence is itself evidence for §8.

\hypertarget{the-measurement-problem-a-structural-discriminator}{\subsection{The measurement problem: a structural discriminator}\label{the-measurement-problem-a-structural-discriminator}}

(From the SM-track extension; included for completeness of the breadth claim.) Measurement-selection --- ``which outcome
becomes the fact'' --- is proven an \textbf{irreducible strict extension} of quantum mechanics in the calculus: records are
\(\Sigma_f\)-definable, selection is not, by an explicit non-factorization witness. The constructive payoff is a
\emph{no-fitting} discriminator: any future single-outcome theory must either modify the defect ledger \(D\)
(objective-collapse family) or only the readout \(\Sigma_f\) (many-worlds family); a community-accepted single-outcome
theory touching \emph{neither} would falsify the claim. This stakes a structural prediction about \emph{where any future theory
must land}, before such a theory exists.

\hypertarget{the-pattern}{\subsection{The pattern}\label{the-pattern}}

\begin{longtable}[]{@{}llll@{}}
\toprule
\begin{minipage}[b]{0.17\columnwidth}\raggedright
puzzle\strut
\end{minipage} & \begin{minipage}[b]{0.15\columnwidth}\raggedright
law\strut
\end{minipage} & \begin{minipage}[b]{0.32\columnwidth}\raggedright
verdict\strut
\end{minipage} & \begin{minipage}[b]{0.26\columnwidth}\raggedright
control\strut
\end{minipage}\tabularnewline
\midrule
\endhead
\begin{minipage}[t]{0.17\columnwidth}\raggedright
Bell nonlocality\strut
\end{minipage} & \begin{minipage}[t]{0.15\columnwidth}\raggedright
F49 + F37\strut
\end{minipage} & \begin{minipage}[t]{0.32\columnwidth}\raggedright
nonlocal residual recovered; carrier \(\neq\) hidden variable (forbidden rule)\strut
\end{minipage} & \begin{minipage}[t]{0.26\columnwidth}\raggedright
classical mixture: locally explainable, LHV table exhibited\strut
\end{minipage}\tabularnewline
\begin{minipage}[t]{0.17\columnwidth}\raggedright
BH information\strut
\end{minipage} & \begin{minipage}[t]{0.15\columnwidth}\raggedright
F34\strut
\end{minipage} & \begin{minipage}[t]{0.32\columnwidth}\raggedright
loss gauge-only at full boundary; real at partial readouts\strut
\end{minipage} & \begin{minipage}[t]{0.26\columnwidth}\raggedright
exact complement-unitary witness\strut
\end{minipage}\tabularnewline
\begin{minipage}[t]{0.17\columnwidth}\raggedright
cosmological constant\strut
\end{minipage} & \begin{minipage}[t]{0.15\columnwidth}\raggedright
F50 (+F47/F26)\strut
\end{minipage} & \begin{minipage}[t]{0.32\columnwidth}\raggedright
bounded-moduli; no in-layer value-law\strut
\end{minipage} & \begin{minipage}[t]{0.26\columnwidth}\raggedright
\(\kappa\): closure-constrained boundary\strut
\end{minipage}\tabularnewline
\begin{minipage}[t]{0.17\columnwidth}\raggedright
measurement\strut
\end{minipage} & \begin{minipage}[t]{0.15\columnwidth}\raggedright
non-factorization\strut
\end{minipage} & \begin{minipage}[t]{0.32\columnwidth}\raggedright
irreducible strict extension; \(D\)-vs-\(\Sigma_f\) discriminator\strut
\end{minipage} & \begin{minipage}[t]{0.26\columnwidth}\raggedright
identical-diagonal witness pair\strut
\end{minipage}\tabularnewline
\bottomrule
\end{longtable}

One grammar; six foundational puzzles (counting §6's three); every verdict computed, controlled, and graded.
 \hypertarget{cross-track-ties-and-the-one-grammar-thesis}{\section{Cross-track ties, and the one-grammar thesis}\label{cross-track-ties-and-the-one-grammar-thesis}}

The per-track results stand on their own. But the calculus offers something that neither track shows alone and that, to
our knowledge, no other framework in either domain produces: \textbf{the same structural object turning out to be two
different physical phenomena}, on substrates that share no physics. These cross-track ties are not analogies or shared
tooling --- they are single objects in one grammar, read on two substrates that were explicitly guarded against each
other. We catalog them first (§8.2), argue why they are distinctive rather than coincidental (§8.3), and then draw the
thesis they support (§8.4).

\hypertarget{why-the-ties-are-non-trivial-the-anti-contamination-guard}{\subsection{Why the ties are non-trivial: the anti-contamination guard}\label{why-the-ties-are-non-trivial-the-anti-contamination-guard}}

The evidential weight of ``one object, two phenomena'' collapses if the two constructions were allowed to copy each
other. They were not. An \textbf{anti-contamination protocol} was declared at the outset: the SM substrate is a
\emph{selection/measure} layer (\(P_2\)-dominant --- a candidate space, constraints, a measure); the QM--GR substrate is a
\emph{co-sourcing common refinement} (\(P_4\)/F51 --- one field, two readouts). Build validators on the SM track fail on any
field-carrier token; steps that pattern-matched one substrate onto the other were rejected during construction
(Section 3.4). So when a single object surfaces on both tracks, it is because the \emph{law} is layer-agnostic --- not because
one construction borrowed from the other. Every tie below survived that guard.

\hypertarget{the-cross-track-ties}{\subsection{The cross-track ties}\label{the-cross-track-ties}}

\begin{quote}
\textbf{Each row is one object, computed on both substrates.} Read across: the leftmost column is a single construct of the
calculus; the next two columns are what it \emph{is} on each track.
\end{quote}

\begin{longtable}[]{@{}lll@{}}
\toprule
\begin{minipage}[b]{0.23\columnwidth}\raggedright
one calculus object\strut
\end{minipage} & \begin{minipage}[b]{0.35\columnwidth}\raggedright
on the SM substrate (selection layer)\strut
\end{minipage} & \begin{minipage}[b]{0.35\columnwidth}\raggedright
on the QM--GR substrate (common refinement)\strut
\end{minipage}\tabularnewline
\midrule
\endhead
\begin{minipage}[t]{0.23\columnwidth}\raggedright
\textbf{factorization defect} \(\Delta_{\mathrm{fact}}\)\strut
\end{minipage} & \begin{minipage}[t]{0.35\columnwidth}\raggedright
the \(X/Y\) colored coset --- proton decay, monopole, color-separation, GUT non-relabeling, one object (§4.4)\strut
\end{minipage} & \begin{minipage}[t]{0.35\columnwidth}\raggedright
the ladder defect (\(8\) witnesses) + the route-mismatch \(\Delta_3^{\mathrm{comp}}\) --- \emph{non-renormalizability} (§5.1--5.2)\strut
\end{minipage}\tabularnewline
\begin{minipage}[t]{0.23\columnwidth}\raggedright
\textbf{the record / memory rung}\strut
\end{minipage} & \begin{minipage}[t]{0.35\columnwidth}\raggedright
grounds clean-separation, \emph{forbids proton decay \& monopoles} (Pred. 3); hosts the measurement discriminator\strut
\end{minipage} & \begin{minipage}[t]{0.35\columnwidth}\raggedright
grounds the common carrier; makes the gravitational channel a record channel --- \emph{forbids BMV entanglement} (Pred. 2)\strut
\end{minipage}\tabularnewline
\begin{minipage}[t]{0.23\columnwidth}\raggedright
\textbf{F51 common refinement}\strut
\end{minipage} & \begin{minipage}[t]{0.35\columnwidth}\raggedright
the GUT parent → \(\sin^2\theta_W = 3/8\), \(k_Y = 5/3\) (trace ratios)\strut
\end{minipage} & \begin{minipage}[t]{0.35\columnwidth}\raggedright
the parent \(L\) → fork-over-ladder; forces the monogamy combiner\strut
\end{minipage}\tabularnewline
\begin{minipage}[t]{0.23\columnwidth}\raggedright
\textbf{GROUND landing mode}\strut
\end{minipage} & \begin{minipage}[t]{0.35\columnwidth}\raggedright
clean-separation \(\leftarrow\) record-stability\strut
\end{minipage} & \begin{minipage}[t]{0.35\columnwidth}\raggedright
common carrier \(\leftarrow\) semiclassical co-sourcing\strut
\end{minipage}\tabularnewline
\begin{minipage}[t]{0.23\columnwidth}\raggedright
\textbf{F37 complementarity} (non-joint-access)\strut
\end{minipage} & \begin{minipage}[t]{0.35\columnwidth}\raggedright
(the access-quotient discipline of the selection layer)\strut
\end{minipage} & \begin{minipage}[t]{0.35\columnwidth}\raggedright
the premise controls (no joint quotient for non-commuting accesses); reconciles the carrier with Bell (§7.1)\strut
\end{minipage}\tabularnewline
\begin{minipage}[t]{0.23\columnwidth}\raggedright
\textbf{F47/F26 demarcation} (no value-law on contingent selection)\strut
\end{minipage} & \begin{minipage}[t]{0.35\columnwidth}\raggedright
no value-law on the EW scale, \(N_{\mathrm{gen}}\), content\strut
\end{minipage} & \begin{minipage}[t]{0.35\columnwidth}\raggedright
no in-layer \(\Lambda\) value-law (F50 bounded-moduli, §7.3)\strut
\end{minipage}\tabularnewline
\bottomrule
\end{longtable}

Three of these are physical statements, not bookkeeping:

\begin{itemize}
\item
  \textbf{One defect is both proton physics and non-renormalizability.} \(\Delta_{\mathrm{fact}}\) --- the failure of one
  readout to factor through another --- is, on the SM substrate, the \(X/Y\) coset whose \emph{absence} makes the proton stable;
  on the QM--GR substrate, it is the obstruction that makes quantizing gravity fail. The calculus says these two
  century-old difficulties are the \emph{same kind of structural fact} about layered descriptions, surfacing on different
  substrates.
\item
  \textbf{One rung underlies three of the paper's headline results.} Record/memory stability --- the requirement that a
  universe be able to hold stable records --- grounds the SM gauge selection (§4.5), forbids proton decay (Prediction 3),
  \emph{and} forces the gravitational channel to be a record channel that decoheres without entangling (Prediction 2). A
  single structural commitment, stated once, has consequences in particle physics \emph{and} in tabletop gravity. This is
  the most consequential tie: it predicts that \emph{whatever forbids the proton from decaying is the same structure that
  prevents gravity from entangling masses} --- a connection no separate treatment of the two would suggest.
\item
  \textbf{One refinement law lands a gauge angle and resolves quantum gravity's architecture.} F51 (unification as common
  refinement) builds the GUT parent that yields \(\sin^2\theta_W = 3/8\) on one substrate and the QM--GR parent that
  selects the fork on the other --- the same commuting-square construction, two unrelated payoffs.
\end{itemize}

\hypertarget{why-this-is-distinctive-and-not-coincidence}{\subsection{Why this is distinctive --- and not coincidence}\label{why-this-is-distinctive-and-not-coincidence}}

Two reasons these ties are evidence rather than wordplay. First, \textbf{the anti-contamination guard} (§8.1): the objects
recur across substrates that were prevented from sharing structure, so the recurrence is forced by the laws. Second,
\textbf{the ties do work} --- each is load-bearing in a \emph{result}, not decorative: the shared record rung is what lets
Predictions 2 and 3 use one record constraint without collapsing their distinct physical claims; the shared \(\Delta_{\mathrm{fact}}\) is what lets the
non-renormalizability typing and the proton-stability spine be checked by one piece of frozen machinery. A coincidence
of notation could do neither.

This is, we contend, a distinctive feature of a layer-agnostic calculus, and one not easily reproduced by
domain-specific methods: a quantum-gravity framework has no reason to predict proton stability, and a
grand-unified model has no reason to constrain gravitationally-induced entanglement, but a calculus whose objects are
substrate-independent \emph{forces} such links and states them as predictions (§6).

\hypertarget{the-one-grammar-thesis}{\subsection{The one-grammar thesis}\label{the-one-grammar-thesis}}

The ties support a single claim --- the thesis of the foundations program (visible in the title of its law catalog:
\emph{layer-agnostic} structural laws):

\begin{quote}
\textbf{One frozen grammar --- quotients, closures, descent, factorization defects, common refinements, fiber volumes,
landing modes --- applied to the two most structurally unrelated deep problems in physics, produced foundational-grade
results on both, in the two opposite genres its own laws anticipated in advance, tied them through shared structural
objects, and yielded three falsifiable mainstream-contradicting predictions.}
\end{quote}

Layer-agnosticism cannot be tested on one substrate; it is tested by freezing the grammar and pointing it at problems
that share nothing --- a deliberately conducted \(n=2\) experiment. And the result was not indiscriminate success but
\textbf{differential, predicted success}: the framework's own laws assign the two problems \emph{different expected genres}.
F47/F26 say a \textbf{contingent selection} supports no value-law --- so the SM track should yield demarcations, blindness
theorems, and conditional selections, with its value-law attempts failing or landing only as embedding ratios. A
\textbf{co-sourcing refinement} supports a positive relation between its two readouts --- so the QM--GR track is where an
\(E=mc^2\)-shaped law could land. Both expectations were recorded \emph{before} the corresponding constructions ran (the
value-law push was explicitly parked on the SM track and redirected to QM--GR on exactly this reasoning), and both were
borne out: the SM track delivered the demarcation profile (including the failures of its value-law attempts), the
QM--GR track landed the ledger/shadow-price relation (review-settled). The grammar \textbf{forbids different things in
different places, and its successes and failures fell along its own forbidding} --- the signature of a meta-theory with
content, not a flexible vocabulary. The three falsifiable predictions then arose from the same final move applied to
each track's strongest theorem, under the same symmetric-scrutiny discipline.

\hypertarget{the-honest-bound-on-the-thesis}{\subsection{The honest bound on the thesis}\label{the-honest-bound-on-the-thesis}}

\begin{itemize}
\tightlist
\item
  \(n=2\) is evidence, not proof, of layer-agnosticism; the atlas of foreclosures contains dozens of further edges, each
  a test the thesis could fail.
\item
  Both substrates are finite toys; frame-transfer is unproven on both. The grammar's universality is demonstrated at
  the level of \emph{structure recovered, demarcated, tied, and predicted} --- not of nature derived.
\item
  The construction program was adversarially reviewed at both ends --- two QM--GR bundles settled and the SM core passed
  review --- but has not been independently replicated; the §6 predictions await independent review.
\item
  The ties are structural identities \emph{within the calculus}; we do not claim the SM and quantum gravity are ``the same
  physics'', only that one structural grammar describes both and links them --- and that the links are falsifiable
  (Predictions 2--3 are the shared record rung made testable).
\end{itemize}

\begin{figure}[tbp]
  \centering
  \includegraphics[width=\linewidth]{figures/fig_f1_one_grammar.png}
  \caption{\textbf{The paper's thesis in one panel.} \emph{Center: the grammar (six primitives + F-laws, one box). Two arrows out, through an ``anti-contamination guard'' wall: left to the SM substrate (selection layer) yielding the demarcation stack (§4) and Prediction 3; right to the QM--GR substrate (common refinement) yielding the unification stack (§5) and Predictions 1--2. Beneath each arrow, the prior expectation (F47/F26: ``no value-law here''; co-sourcing: ``a law can land here'') with a timestamp glyph indicating the expectation preceded the result. Connecting the two arms: the cross-track ties of §8.2 drawn as rungs between the substrates --- labeled \(\Delta_{\mathrm{fact}}\), the record rung, F51 --- with the record rung highlighted as the spine of Predictions 2 and 3.}}
  \label{fig:f1}
\end{figure}
 \hypertarget{scope-limits-and-falsifiability}{\section{Scope, limits, and falsifiability}\label{scope-limits-and-falsifiability}}

This section collects every bound in one place, then gives the claim-by-claim falsifiability table. These are scope
statements, stated flat; none retracts a result.

\hypertarget{the-bounds}{\subsection{The bounds}\label{the-bounds}}

\begin{enumerate}
\def\labelenumi{\arabic{enumi}.}
\tightlist
\item
  \textbf{Finite toys; no frame-transfer.} Every computation lives on a declared finite carrier. The carriers \emph{predict}
  and are \emph{falsifiable}, but nothing here is proven about nature. The three predictions of §6 are exactly that ---
  predictions; their status in nature is for experiment and for richer constructions to decide.
\item
  \textbf{Conditional landings are conditional.} Every GROUND/SELECT result rests on a named source: clean-separation on
  record-stability (``why must a universe bear records?'' is the review-judged top of that tower); the QM--GR fork on
  the co-sourcing warrant (which reaches matter-sourced geometry, not the free gravitational sector). The conditions
  are declared, audited, and carried with the results.
\item
  \textbf{Enumeration-strength where stated.} The record-stability implication is \(0/11{,}990\) with bounded outside probes
  --- not yet an all-structures theorem; the open lemma (L60\(\to\)L64: mass-closure charge-linking forbids a second
  neutral record channel in leak-positive substrate closers) is precisely specified and would upgrade Prediction 3
  from enumeration-strength to a structural law. Similar precisely-typed upgrades are open for: exact RT saturation in
  the large-bond limit; the continuum (F43) limit of the discrete Einstein condition; uniqueness of the co-sourcing
  grounding; the all-co-readout form of the monogamy forcing.
\item
  \textbf{\(n=2\) on universality.} Section 8's thesis has exactly two substrates behind it.
\item
  \textbf{Adversarially reviewed, not independently replicated.} Two results settled under adversarial review cycles; no
  independent group has yet re-run the constructions (Appendix B is written to make that easy).
\item
  \textbf{What is \emph{not} claimed anywhere:} a derivation of any measured constant; a solution of quantum gravity; a proof of
  physical proton stability; a universal exact-six reduction; frame-transfer.
\end{enumerate}

\begin{landscape}

\hypertarget{table-t1-every-headline-claim-its-grade-and-what-kills-it}{\subsection{Table T1 --- every headline claim, its grade, and what kills it}\label{table-t1-every-headline-claim-its-grade-and-what-kills-it}}

\begin{longtable}[]{@{}llll@{}}
\toprule
\begin{minipage}[b]{0.05\columnwidth}\raggedright
\#\strut
\end{minipage} & \begin{minipage}[b]{0.26\columnwidth}\raggedright
claim\strut
\end{minipage} & \begin{minipage}[b]{0.22\columnwidth}\raggedright
grade\strut
\end{minipage} & \begin{minipage}[b]{0.40\columnwidth}\raggedright
what would falsify / defeat it\strut
\end{minipage}\tabularnewline
\midrule
\endhead
\begin{minipage}[t]{0.05\columnwidth}\raggedright
1\strut
\end{minipage} & \begin{minipage}[t]{0.26\columnwidth}\raggedright
\(\sim99.3\%\) chiral-gauge exclusion → \(\{2|3,\ SU(4)\}\)\strut
\end{minipage} & \begin{minipage}[t]{0.22\columnwidth}\raggedright
COMPUTE (unconditional, in-carrier)\strut
\end{minipage} & \begin{minipage}[t]{0.40\columnwidth}\raggedright
a neutral-closure error: a structure passing all gates outside the family (re-run Appendix B.1)\strut
\end{minipage}\tabularnewline
\begin{minipage}[t]{0.05\columnwidth}\raggedright
2\strut
\end{minipage} & \begin{minipage}[t]{0.26\columnwidth}\raggedright
SM gauge algebra selected on clean-separation\strut
\end{minipage} & \begin{minipage}[t]{0.22\columnwidth}\raggedright
SELECT (conditional)\strut
\end{minipage} & \begin{minipage}[t]{0.40\columnwidth}\raggedright
a clean-separated competitor in-window; or the grounding (row 3) failing\strut
\end{minipage}\tabularnewline
\begin{minipage}[t]{0.05\columnwidth}\raggedright
3\strut
\end{minipage} & \begin{minipage}[t]{0.26\columnwidth}\raggedright
clean-separation grounded in record-stability\strut
\end{minipage} & \begin{minipage}[t]{0.22\columnwidth}\raggedright
GROUND, enumeration-strength\strut
\end{minipage} & \begin{minipage}[t]{0.40\columnwidth}\raggedright
\textbf{one} record-stable structure with \(\Delta_{\mathrm{fact}}\neq\varnothing\) (in-carrier or in any lawful extension); refutation of L60→L64 closes the door the other way\strut
\end{minipage}\tabularnewline
\begin{minipage}[t]{0.05\columnwidth}\raggedright
4\strut
\end{minipage} & \begin{minipage}[t]{0.26\columnwidth}\raggedright
\(\sin^2\theta_W=3/8\), \(k_Y=5/3\) as forced embedding ratios\strut
\end{minipage} & \begin{minipage}[t]{0.22\columnwidth}\raggedright
GROUND (minimal-simple source)\strut
\end{minipage} & \begin{minipage}[t]{0.40\columnwidth}\raggedright
showing the trace ratios depend on an undeclared input; the product control (\(3/23\)) failing to discriminate\strut
\end{minipage}\tabularnewline
\begin{minipage}[t]{0.05\columnwidth}\raggedright
5\strut
\end{minipage} & \begin{minipage}[t]{0.26\columnwidth}\raggedright
\(N_{\mathrm{gen}}\)-blindness, all \(N\)\strut
\end{minipage} & \begin{minipage}[t]{0.22\columnwidth}\raggedright
PROVE-BLIND (theorem)\strut
\end{minipage} & \begin{minipage}[t]{0.40\columnwidth}\raggedright
an \(N\)-sensitive gate in the frozen chain for an even-doublet anomaly-free unit\strut
\end{minipage}\tabularnewline
\begin{minipage}[t]{0.05\columnwidth}\raggedright
6\strut
\end{minipage} & \begin{minipage}[t]{0.26\columnwidth}\raggedright
QM--GR fork (no weak-RT ladder)\strut
\end{minipage} & \begin{minipage}[t]{0.22\columnwidth}\raggedright
COMPUTE (conditional on premise)\strut
\end{minipage} & \begin{minipage}[t]{0.40\columnwidth}\raggedright
a lawful factorization \(q_{\mathrm{GR}}=\phi\circ q_{\mathrm{QM}}\) on the carrier; the nested control failing to flip\strut
\end{minipage}\tabularnewline
\begin{minipage}[t]{0.05\columnwidth}\raggedright
7\strut
\end{minipage} & \begin{minipage}[t]{0.26\columnwidth}\raggedright
common-carrier premise grounded in co-sourcing\strut
\end{minipage} & \begin{minipage}[t]{0.22\columnwidth}\raggedright
GROUND (review-settled)\strut
\end{minipage} & \begin{minipage}[t]{0.40\columnwidth}\raggedright
the warrant shown circular (co-sourcing presupposing the carrier); a complementary pair admitting a joint quotient\strut
\end{minipage}\tabularnewline
\begin{minipage}[t]{0.05\columnwidth}\raggedright
8\strut
\end{minipage} & \begin{minipage}[t]{0.26\columnwidth}\raggedright
RT = ledger/shadow-price duality\strut
\end{minipage} & \begin{minipage}[t]{0.22\columnwidth}\raggedright
RECOGNITION (review-settled)\strut
\end{minipage} & \begin{minipage}[t]{0.40\columnwidth}\raggedright
exact (machine-\(\epsilon\)) saturation at finite \(D\) --- the tautology signature; or geometric side tracking the state spectrum under reseeding\strut
\end{minipage}\tabularnewline
\begin{minipage}[t]{0.05\columnwidth}\raggedright
9\strut
\end{minipage} & \begin{minipage}[t]{0.26\columnwidth}\raggedright
\textbf{Prediction 1: entanglement \(\nRightarrow\) geometry}\strut
\end{minipage} & \begin{minipage}[t]{0.22\columnwidth}\raggedright
FORBIDDEN-RULE (toy)\strut
\end{minipage} & \begin{minipage}[t]{0.40\columnwidth}\raggedright
a non-trivial holographic geometry with injective fingerprint map (zero kernel, no two-sided invisible deformation)\strut
\end{minipage}\tabularnewline
\begin{minipage}[t]{0.05\columnwidth}\raggedright
10\strut
\end{minipage} & \begin{minipage}[t]{0.26\columnwidth}\raggedright
\textbf{Prediction 2: gravity does not entangle (BMV-null)}\strut
\end{minipage} & \begin{minipage}[t]{0.22\columnwidth}\raggedright
FORBIDDEN-RULE (toy; analytically verified)\strut
\end{minipage} & \begin{minipage}[t]{0.40\columnwidth}\raggedright
a confirmed BMV-type detection of gravitationally-induced entanglement; or the coherent-mediator control failing to entangle\strut
\end{minipage}\tabularnewline
\begin{minipage}[t]{0.05\columnwidth}\raggedright
11\strut
\end{minipage} & \begin{minipage}[t]{0.26\columnwidth}\raggedright
\textbf{Prediction 3: records forbid proton decay \& monopoles}\strut
\end{minipage} & \begin{minipage}[t]{0.22\columnwidth}\raggedright
FORBIDDEN-RULE / DEMARCATION (enumeration-strength)\strut
\end{minipage} & \begin{minipage}[t]{0.40\columnwidth}\raggedright
an observed proton decay or magnetic monopole; or one record-stable breaking structure\strut
\end{minipage}\tabularnewline
\begin{minipage}[t]{0.05\columnwidth}\raggedright
12\strut
\end{minipage} & \begin{minipage}[t]{0.26\columnwidth}\raggedright
Bell: carrier \(\neq\) local hidden variable\strut
\end{minipage} & \begin{minipage}[t]{0.22\columnwidth}\raggedright
FORBIDDEN-RULE (recovered Bell + reconciliation)\strut
\end{minipage} & \begin{minipage}[t]{0.40\columnwidth}\raggedright
an admissible joint quotient for complementary accesses; an LHV table reproducing the quantum CHSH\strut
\end{minipage}\tabularnewline
\begin{minipage}[t]{0.05\columnwidth}\raggedright
13\strut
\end{minipage} & \begin{minipage}[t]{0.26\columnwidth}\raggedright
BH information: loss is readout-relative\strut
\end{minipage} & \begin{minipage}[t]{0.22\columnwidth}\raggedright
FORBIDDEN-RULE (static toy)\strut
\end{minipage} & \begin{minipage}[t]{0.40\columnwidth}\raggedright
a gauge-invariant interior observable unrecoverable from the full boundary on the full-rank stratum\strut
\end{minipage}\tabularnewline
\begin{minipage}[t]{0.05\columnwidth}\raggedright
14\strut
\end{minipage} & \begin{minipage}[t]{0.26\columnwidth}\raggedright
\(\Lambda\): no in-layer value-law\strut
\end{minipage} & \begin{minipage}[t]{0.22\columnwidth}\raggedright
DEMARCATION (F50 bounded-moduli)\strut
\end{minipage} & \begin{minipage}[t]{0.40\columnwidth}\raggedright
the frozen closure chain selecting a singleton background\strut
\end{minipage}\tabularnewline
\begin{minipage}[t]{0.05\columnwidth}\raggedright
15\strut
\end{minipage} & \begin{minipage}[t]{0.26\columnwidth}\raggedright
one-grammar thesis\strut
\end{minipage} & \begin{minipage}[t]{0.22\columnwidth}\raggedright
meta-claim, \(n=2\)\strut
\end{minipage} & \begin{minipage}[t]{0.40\columnwidth}\raggedright
a third substrate where the grammar's own genre-expectation fails; or evidence of cross-track contamination\strut
\end{minipage}\tabularnewline
\begin{minipage}[t]{0.05\columnwidth}\raggedright
16\strut
\end{minipage} & \begin{minipage}[t]{0.26\columnwidth}\raggedright
measurement: \(D\)-vs-\(\Sigma_f\) discriminator\strut
\end{minipage} & \begin{minipage}[t]{0.22\columnwidth}\raggedright
structural prediction\strut
\end{minipage} & \begin{minipage}[t]{0.40\columnwidth}\raggedright
a community-accepted single-outcome theory touching neither \(D\) nor \(\Sigma_f\)\strut
\end{minipage}\tabularnewline
\bottomrule
\end{longtable}
\end{landscape}

\hypertarget{why-the-honesty-is-load-bearing}{\subsection{Why the honesty is load-bearing}\label{why-the-honesty-is-load-bearing}}

A framework this general invites the Pauli verdict (``not even wrong''). The table above is the rebuttal format: each
claim is either graded as conditional/toy (and says what discharges the condition), or names an observation that would
falsify it. Rows 9, 10, and 11 are the paper's exposure to experiment --- with row 10 (BMV) the nearest-term; rows 3, 6, 7, 8 are
its exposure to construction; row 15 is its exposure to its own program.
 \hypertarget{discussion-and-outlook}{\section{Discussion and outlook}\label{discussion-and-outlook}}

\hypertarget{if-the-predictions-hold-if-they-fail}{\subsection{If the predictions hold; if they fail}\label{if-the-predictions-hold-if-they-fail}}

\textbf{Prediction 1.} If entanglement-invisible geometry is confirmed in richer settings --- higher-rank tensor networks,
genuine AdS/CFT setups, eventually any operational probe of bulk reconstruction --- the consequence is architectural:
reconstruction programs need a second boundary currency beyond entanglement, and ``geometry from entanglement'' becomes
``geometry from entanglement \emph{plus an irreducible remainder}'', with the fork supplying the structural reason. If instead
someone exhibits a non-trivial geometry with an injective fingerprint map, the fork's sharp form fails: the calculus
would have over-claimed the independence of the gravitational mode, and §5.2's architecture --- the paper's center ---
would need revision. The prediction can be tested without experiment: it requires only a construction.

\textbf{Prediction 2 (BMV).} This is the nearest-term test. If a BMV-class experiment detects gravitationally-induced
entanglement between masses, the fork's reading of gravity as a classically-indexed record channel is refuted directly ---
gravity carries quantum phase coherently after all, and §5.2's architecture would need revision. If instead such
experiments find decoherence without entanglement, the prediction holds and the inference ``BMV-positive \(\Rightarrow\)
quantized mediator'' is, at minimum, shown to be premise-dependent. Because the experiments are funded and advancing, the
framework's exposure here is on a decade rather than a cosmological timescale.

\textbf{Prediction 3.} If proton decay is observed at Hyper-Kamiokande or a monopole is confirmed, the record-stability link
is severed and the grounding of §4.5 collapses --- taking with it the conditional selection of the SM gauge structure as
this paper grounds it. That is the intended exposure: the framework's deepest SM result is hostage to a running
experiment. Conversely, every year of proton-decay null results is consistent with a framework that --- uniquely, to our
knowledge --- predicts \emph{exact} stability for a structural reason tied to the existence of records.

\hypertarget{nearest-upgrades-precisely-typed-open-problems}{\subsection{Nearest upgrades (precisely-typed open problems)}\label{nearest-upgrades-precisely-typed-open-problems}}

The program's open problems are not vague: each is a named lemma or audit with a stated success criterion. (i) The
\textbf{L60\(\to\)L64} lemma --- mass-closure charge-linking forbids a second neutral record channel --- upgrades Prediction 3 to
an all-structures law. (ii) The \textbf{F43 continuum audit} --- tower-coherence, cofinality, residual-vanishing, completion ---
either licenses the continuum Einstein limit of the discrete consistency condition or names the exact obstruction.
(iii) \textbf{Exact RT saturation} in the large-bond limit (the \(0.729\to0.904\to0.941\) trend made a theorem with a rate).
(iv) \textbf{Uniqueness of the grounding}: whether semiclassical co-sourcing is the only admissible warrant for the common
carrier (warranted \(\to\) forced). (v) A \textbf{third substrate} for the universality thesis --- the atlas's remaining
foreclosure edges (black-hole thermodynamics and singularities are the natural candidates) --- turning \(n=2\) into \(n=3\)
with the genre-expectation again stated in advance.

\hypertarget{what-kind-of-theory-is-this}{\subsection{What kind of theory is this?}\label{what-kind-of-theory-is-this}}

It is worth being precise about the genre, because misreading it produces both over- and under-claims. The calculus is
not a dynamical theory and competes with none: it has no Lagrangian, predicts no cross-sections, and cannot replace
quantum field theory or general relativity. It is a \textbf{meta-theory of inter-layer structure} --- of what can and cannot
descend, compose, refine, and close between descriptions. Its natural outputs are demarcations (``this is observed-input,
provably''), unifications (``these known results are one law''), forbidden rules (``this configuration cannot occur''), and
--- as Section 6 shows --- falsifiable predictions where its structural laws collide with the commitments of dynamical
programs. The three collisions exhibited here (with it-from-qubit, with the BMV-positive expectation, and with grand unification) are, we suggest, the
beginning of a useful division of labor: dynamical theories propose mechanisms; a structural calculus rules on which
mechanisms are architecturally possible --- and is refutable when its rulings are wrong.

\hypertarget{an-invitation}{\subsection{An invitation}\label{an-invitation}}

Appendix B contains everything needed to re-run every number in this paper deterministically, and Table T1 names what
would falsify each claim. The fastest meaningful attacks: construct the injective-fingerprint geometry (defeats Prediction 1);
exhibit a record-stable breaking structure in any lawful extension of the carrier (defeats Prediction 3 at
enumeration level); prove or refute L60\(\to\)L64 (settles Prediction 3's theorem status); or run the grammar on a third
substrate and watch whether its genre-expectation fails (attacks the thesis of §8). We would consider any of these
outcomes --- confirmation or refutation --- a success of the method.
 \hypertarget{conclusion}{\section{Conclusion}\label{conclusion}}

We applied one frozen, layer-agnostic structural calculus --- six primitive inter-layer roles, a catalog of
substrate-independent laws, a mechanical audit discipline --- to physics' two deepest and most unrelated foreclosures,
under an anti-contamination protocol that prevented the constructions from informing each other.

The calculus produced, on finite reproducible carriers: for the \textbf{Standard Model}, a demarcation --- an unconditional
exclusion of \(\sim99.3\%\) of chiral gauge structures, a conditional selection of
\(\mathfrak{su}(2)\oplus\mathfrak{su}(3)\oplus\mathfrak{u}(1)\) grounded in record-stability, the non-circular recovery of
\(\sin^2\theta_W = 3/8\) and \(k_Y = 5/3\), and proof-grade blindness theorems certifying the generation count and matter
content as observed-input; for \textbf{QM--GR}, a unification --- the non-composition located as a non-commuting-completions
defect, resolved by a unique common-refinement fork, with Ryu--Takayanagi recognized (and review-checked) as a
special case of one ledger/shadow-price law. The same machinery then classified Bell nonlocality, black-hole
information, and the cosmological constant --- and the two result-genres landed exactly where the calculus's own laws,
on record in advance, said they would.

The principal results are three forced, falsifiable, mainstream-contradicting predictions. \textbf{Entanglement does not
determine geometry}: the entanglement\(\to\)geometry map carries a robust kernel (dimension 5 on the test carrier, with
two-sided finite invisible deformations), so complete entanglement-based bulk reconstruction is impossible --- refutable
by a single injective counter-geometry. \textbf{Gravity does not mediate entanglement}: the same fork forces the
gravitational channel to be a classically-indexed record channel, so BMV-class experiments are predicted null
(decoherence without entanglement; verified against closed-form quantum mechanics) --- refutable on tabletop timescales,
the nearest-term exposure of the three. \textbf{Records forbid proton decay and monopoles}: record-bearing structures lie
provably in the branch without the \(X/Y\) coset (\(0/11{,}990\), non-circularly), so the proton of a record-bearing
universe is exactly stable --- refutable by one decay event or one monopole. A single record/memory layer ties the
predictions at the grammar level without identifying the two physical tracks.

Everything is bounded as stated: finite toys, conditional landings, enumeration-strength where flagged, \(n=2\) on
universality, no frame-transfer claimed --- and every claim carries a named falsifier. The contribution, in one sentence:
\emph{a structural meta-theory of emergence, applied under adversarial self-correction, produces foundational-grade
results on both of physics' deepest problems and yields refutable, mainstream-contradicting predictions on the outcome.}
The predictions are now the community's to test.
 
\appendix
\hypertarget{the-calculus-formally-working-subset}{\section{the calculus, formally (working subset)}\label{the-calculus-formally-working-subset}}

This appendix states, at usable precision, every formal object the body relies on. Full development: Foundations II--IV~\citep{Tsiokos_FoundationsII,Tsiokos_FoundationsIII,Tsiokos_FoundationsIV}.

\hypertarget{descent-and-obstruction}{\subsection{Descent and obstruction}\label{descent-and-obstruction}}

For a surjection \(q:H\to Q\) and readout \(t:H\to T\):

\[
t \text{ descends} \iff \mathcal O_q(t) = \{(h,h'): q(h)=q(h') \wedge t(h)\neq t(h')\} = \varnothing
\iff \exists\,\bar t:Q\to T,\ \bar t\circ q = t .
\]

Both directions are finite checks on \(H\times H\) (choose representatives; well-definedness is exactly
fiber-constancy).

\hypertarget{strict-extension-and-the-factorization-defect-fiii-thms-1213}{\subsection{Strict extension and the factorization defect (FIII Thms 12--13)}\label{strict-extension-and-the-factorization-defect-fiii-thms-1213}}

For maps \(\pi_0:S\to O_0\), \(\pi_1:S\to O_1\) on a finite carrier:

\[
\Delta_{\mathrm{fact}}(\pi_0,\pi_1) = \{(s,s'): \pi_0(s)=\pi_0(s') \wedge \pi_1(s)\neq\pi_1(s')\},
\qquad
\Delta_{\mathrm{fact}}\neq\varnothing \iff \nexists\,\phi:\ \pi_1 = \phi\circ\pi_0 .
\]

Used as: the \(X/Y\) coset (§4: witnesses of the clean-separation defect); GUT\(\leftrightarrow\)SM non-relabeling
(\(\Delta_{\mathrm{fact}}(\mathrm{SM},\mathrm{GUT})=15\), reverse \(=0\)); the QM--GR ladder defect (\(8\) witnesses).

\hypertarget{non-commuting-completions-fiii-thms-910}{\subsection{Non-commuting completions (FIII Thms 9--10)}\label{non-commuting-completions-fiii-thms-910}}

There exist idempotents \(E_1,E_2\) on a finite carrier with \(E_1E_2\neq E_2E_1\); the route-mismatch defect
\(\Delta_3^{\mathrm{comp}}(E_1,E_2)=\{c: E_1E_2(c)\neq E_2E_1(c)\}\) is non-empty iff they fail to commute. Minimal
witness: \(C=\mathcal P(\{a,b,c\})\), \(E_1(S)=S\cup\{b\}\) if \(a\in S\) else \(S\); \(E_2(S)=S\cup\{c\}\) if \(b\in S\) else \(S\);
then \(E_1E_2(\{a\})=\{a,b\}\neq\{a,b,c\}=E_2E_1(\{a\})\). Reading: route mismatch is a real typed defect even when each
completion is individually exact --- the calculus's typing of non-renormalizability (§5.1).

\hypertarget{the-six-primitives-and-the-no-algebra-theorem}{\subsection{The six primitives and the no-algebra theorem}\label{the-six-primitives-and-the-no-algebra-theorem}}

\(\mathbb P=\{P_1,\dots,P_6\}\) (descent, representability, route mismatch, refinement, packaging, audit) are role
labels for typed judgments. Theorem-grade exclusions (FIII): no total binary operation \(*:\mathbb P^2\to\mathbb P\)
decodes judgment status (no total six-symbol algebra); no universal claim that all structure decomposes into the six.
Directed-cell notation \(P_i \leftarrow P_j\) names a typed, individually-audited judgment, never a product.

\hypertarget{the-f-laws-used-in-this-paper-one-line-formal-cores}{\subsection{The F-laws used in this paper (one-line formal cores)}\label{the-f-laws-used-in-this-paper-one-line-formal-cores}}

\begin{longtable}[]{@{}lll@{}}
\toprule
\begin{minipage}[b]{0.23\columnwidth}\raggedright
law\strut
\end{minipage} & \begin{minipage}[b]{0.51\columnwidth}\raggedright
statement (core)\strut
\end{minipage} & \begin{minipage}[b]{0.19\columnwidth}\raggedright
used in\strut
\end{minipage}\tabularnewline
\midrule
\endhead
\begin{minipage}[t]{0.23\columnwidth}\raggedright
\textbf{F23} probability\strut
\end{minipage} & \begin{minipage}[t]{0.51\columnwidth}\raggedright
a readout unresolved on the fiber of \(q\) admits a \emph{stable} weight \(\Pr_{q,r}(a,v)=m(\lambda(a))(\{h: q(h)=a, r(h)=v\})\)\strut
\end{minipage} & \begin{minipage}[t]{0.19\columnwidth}\raggedright
§5.4 (Born ledger)\strut
\end{minipage}\tabularnewline
\begin{minipage}[t]{0.23\columnwidth}\raggedright
\textbf{F26/F47} contingency \& fine-tuning\strut
\end{minipage} & \begin{minipage}[t]{0.51\columnwidth}\raggedright
no value-law lands on a contingent (moduli) selection; fine-tuning = small selector region\strut
\end{minipage} & \begin{minipage}[t]{0.19\columnwidth}\raggedright
§4.6, §7.3, §8.4\strut
\end{minipage}\tabularnewline
\begin{minipage}[t]{0.23\columnwidth}\raggedright
\textbf{F27} conservation\strut
\end{minipage} & \begin{minipage}[t]{0.51\columnwidth}\raggedright
a charge is conserved iff it descends through the orbit quotient (obstruction \(=0\))\strut
\end{minipage} & \begin{minipage}[t]{0.19\columnwidth}\raggedright
§4.4, §6.3 (baryon number)\strut
\end{minipage}\tabularnewline
\begin{minipage}[t]{0.23\columnwidth}\raggedright
\textbf{F34} information loss\strut
\end{minipage} & \begin{minipage}[t]{0.51\columnwidth}\raggedright
\(\mathrm{Recoverable}(\sigma,\tau,\rho) \iff \mathcal O_\rho=\{(h,h'):\rho(\tau h)=\rho(\tau h') \wedge \sigma(h)\neq\sigma(h')\}=\varnothing\); loss legitimate iff later identifications are invisible to \(\sigma\)\strut
\end{minipage} & \begin{minipage}[t]{0.19\columnwidth}\raggedright
§7.2\strut
\end{minipage}\tabularnewline
\begin{minipage}[t]{0.23\columnwidth}\raggedright
\textbf{F37} complementarity\strut
\end{minipage} & \begin{minipage}[t]{0.51\columnwidth}\raggedright
\(q_A \perp q_B \iff \nexists (J,j,a,b):\ a\circ j=q_A \wedge b\circ j=q_B\) (no admissible joint quotient)\strut
\end{minipage} & \begin{minipage}[t]{0.19\columnwidth}\raggedright
§5.3, §7.1\strut
\end{minipage}\tabularnewline
\begin{minipage}[t]{0.23\columnwidth}\raggedright
\textbf{F39} entropy\strut
\end{minipage} & \begin{minipage}[t]{0.51\columnwidth}\raggedright
\(E(q)\) = fiber volume; chain rule \(E(q_C)=E(q_F)\odot m\); monotone under coarsening\strut
\end{minipage} & \begin{minipage}[t]{0.19\columnwidth}\raggedright
§5.4 (area ledger)\strut
\end{minipage}\tabularnewline
\begin{minipage}[t]{0.23\columnwidth}\raggedright
\textbf{F48} gluing/topology\strut
\end{minipage} & \begin{minipage}[t]{0.51\columnwidth}\raggedright
global obstruction as non-trivial gluing class (monopole obstruction \(=0\) in the clean branch)\strut
\end{minipage} & \begin{minipage}[t]{0.19\columnwidth}\raggedright
§4.4, §6.3\strut
\end{minipage}\tabularnewline
\begin{minipage}[t]{0.23\columnwidth}\raggedright
\textbf{F49} common source\strut
\end{minipage} & \begin{minipage}[t]{0.51\columnwidth}\raggedright
\(\mathrm{LocallyExplainable}(\kappa) \iff \mathrm{CommonSource}(\sigma) \vee \mathrm{InterfaceFactorization}(\iota)\); else statused nonlocal residual\strut
\end{minipage} & \begin{minipage}[t]{0.19\columnwidth}\raggedright
§7.1\strut
\end{minipage}\tabularnewline
\begin{minipage}[t]{0.23\columnwidth}\raggedright
\textbf{F50} background selection\strut
\end{minipage} & \begin{minipage}[t]{0.51\columnwidth}\raggedright
descended background \(b\): \textbf{necessary} (\(b\) constant) xor \textbf{moduli} (\(\exists m,m': b(m)\neq b(m')\)); \textbf{vacuum} iff \(V=\{m_*\}\)\strut
\end{minipage} & \begin{minipage}[t]{0.19\columnwidth}\raggedright
§7.3\strut
\end{minipage}\tabularnewline
\begin{minipage}[t]{0.23\columnwidth}\raggedright
\textbf{F51} unification\strut
\end{minipage} & \begin{minipage}[t]{0.51\columnwidth}\raggedright
parent \(Q_U\) with commuting squares \(\pi_A\circ q_U = q_A\circ\iota_A\), \(\pi_B\circ q_U = q_B\circ\iota_B\) + status-compatible claim lifts\strut
\end{minipage} & \begin{minipage}[t]{0.19\columnwidth}\raggedright
§4.4, §5.2, §5.4\strut
\end{minipage}\tabularnewline
\bottomrule
\end{longtable}

\hypertarget{landing-modes-and-audit-gates-operational-definitions}{\subsection{Landing modes and audit gates (operational definitions)}\label{landing-modes-and-audit-gates-operational-definitions}}

\textbf{Modes.} COMPUTE (value/verdict from the construction); SELECT (choice on declared conditions); PROVE-BLIND
(invariance theorem ⇒ observed-input); GROUND (condition supplied as down-shadow of a named higher source; carries
source-warrant + theorem-strength caveats); RECOGNITION (known result reproduced as a law-instance, non-circularly);
FORBIDDEN-RULE/DEMARCATION (computed exclusion/classification with can-fail control).

\textbf{Gates} (all mechanical; validators enforce): \emph{anti-tautology} --- target not definitionally equal to its evidence;
\emph{anti-circularity} --- each hypothesis independently satisfiable by a conclusion-violating structure (exhibited);
\emph{frozen-machinery} --- imports SHA-256-pinned, no retuning; \emph{can-fail control} --- a configuration on which the claim fails;
\emph{no-overclaim} --- artifact text checked against a blocklist; \emph{anti-contamination} (cross-track) --- substrate-foreign
tokens fail the build.

\hypertarget{the-key-derived-quantities-of-the-body}{\subsection{The key derived quantities of the body}\label{the-key-derived-quantities-of-the-body}}

\begin{itemize}
\tightlist
\item
  \textbf{Embedding ratios} (§4.4): over one SM generation,
  \(\operatorname{Tr}T_3^2 = 3\!\cdot\!2\!\cdot\!\tfrac14 + 2\!\cdot\!\tfrac14\!\cdot\!2 = 2\) (quark + lepton doublets),
  \(\operatorname{Tr}Q^2 = \tfrac83 + \tfrac23 + 2 = \tfrac{16}{3}\), hence \(\sin^2\theta_W = 3/8\); hypercharge
  normalization \(k_Y = 5/3\) from \(\operatorname{Tr}Y^2 = 5/6\). Control: a non-simple product parent yields \(3/23\).
\item
  \textbf{RT/shadow price} (§5.4): for entanglement flow with capacities \(b\), strong duality gives
  \(\max(\text{flow}) = \min(\text{cut})\) and the cut value is the dual optimum --- the shadow price
  \(\lambda = \partial \Phi^*/\partial b\) of the capacity constraints; ``area'' is that price.
\item
  \textbf{Monogamy} (§5.4, §6.1): \(I_3 \le 0\) for min-cut-realizable entropy vectors; GHZ\(_4\) has \(I_3 = +\log 2\).
\item
  \textbf{CHSH} (§7.1): local bound \(\max_{\lambda\in\{\pm1\}^4} |E_{00}+E_{01}+E_{10}-E_{11}| = 2\) (16-strategy
  enumeration); Bell state with \(a_0=Z, a_1=X, b_{0,1}=(Z\pm X)/\sqrt2\) gives \(2\sqrt2\).
\item
  \textbf{Prediction kernels} (§6): \(J = \partial\mathcal E/\partial w\) (central differences, non-degenerate base);
  \(\dim\ker J = 5\); clean-shadow edges = \(\{e: \mathcal E(w \pm \delta\,\hat e) = \mathcal E(w)\ \text{exactly},\
  \delta \in [0.05, 0.15]\}\), count \(2\) per seed.
\end{itemize}
 \hypertarget{toy-constructions-reproducibility}{\section{toy constructions \& reproducibility}\label{toy-constructions-reproducibility}}

Every number in the body regenerates deterministically from the artifact repository accompanying this paper
(\nolinkurl{physics_atlas/}, tracks \nolinkurl{thread_cluster_a/} and \nolinkurl{thread_qm_gr/}). Each step directory contains the build script (the
construction), a machine-readable schema (the claimed numbers), and a validator \nolinkurl{run_stepNN.py} that \textbf{recomputes} the
headline quantities from frozen inputs and fails on mismatch, on anti-smuggling-gate violations, and on overclaim
phrases. Imports between steps are SHA-256-pinned. Python 3 + NumPy suffice; every validator runs offline.

\begin{landscape}

\hypertarget{the-sm-track-selection-layer-key-artifacts}{\subsection{The SM track (selection layer) --- key artifacts}\label{the-sm-track-selection-layer-key-artifacts}}

\begin{longtable}[]{@{}lll@{}}
\toprule
\begin{minipage}[b]{0.21\columnwidth}\raggedright
result (§)\strut
\end{minipage} & \begin{minipage}[b]{0.38\columnwidth}\raggedright
artifact dir (under \nolinkurl{thread_cluster_a/steps/})\strut
\end{minipage} & \begin{minipage}[b]{0.34\columnwidth}\raggedright
headline numbers the validator re-derives\strut
\end{minipage}\tabularnewline
\midrule
\endhead
\begin{minipage}[t]{0.21\columnwidth}\raggedright
neutral carrier + token-blind selection (§4.2)\strut
\end{minipage} & \begin{minipage}[t]{0.38\columnwidth}\raggedright
\nolinkurl{step8_structural_token_blind_selection_*}\strut
\end{minipage} & \begin{minipage}[t]{0.34\columnwidth}\raggedright
\(8640\)-world neutral product; structural pruning \(\to 468\); realized analog rank 2 in a 6-way tie (not unique)\strut
\end{minipage}\tabularnewline
\begin{minipage}[t]{0.21\columnwidth}\raggedright
corrected closure / exclusion (§4.2)\strut
\end{minipage} & \begin{minipage}[t]{0.38\columnwidth}\raggedright
\nolinkurl{step28–42} family (corrected chirality, mass-closure)\strut
\end{minipage} & \begin{minipage}[t]{0.34\columnwidth}\raggedright
\(11{,}990\) genuinely-chiral structures \(\to \{2|3, SU(4)\}\)\strut
\end{minipage}\tabularnewline
\begin{minipage}[t]{0.21\columnwidth}\raggedright
clean separation = \(\Delta_{\mathrm{fact}}\); \(X/Y\) coset (§4.3--4.4)\strut
\end{minipage} & \begin{minipage}[t]{0.38\columnwidth}\raggedright
\nolinkurl{step41_mode_b_factorization_defect_clean_separation_*}\strut
\end{minipage} & \begin{minipage}[t]{0.34\columnwidth}\raggedright
defect witnesses = \(X/Y\) coset; \(\Delta_{\mathrm{fact}}(\mathrm{SM},\mathrm{GUT})=15\), reverse \(0\)\strut
\end{minipage}\tabularnewline
\begin{minipage}[t]{0.21\columnwidth}\raggedright
record-stability grounding (§4.5)\strut
\end{minipage} & \begin{minipage}[t]{0.38\columnwidth}\raggedright
\nolinkurl{step57/58/59_mode_b_record_stability_*}\strut
\end{minipage} & \begin{minipage}[t]{0.34\columnwidth}\raggedright
\(0\) record-stable \(\wedge\) breaking out of \(11{,}990\); substrate-only \(60\); capacity-only \(311\)\strut
\end{minipage}\tabularnewline
\begin{minipage}[t]{0.21\columnwidth}\raggedright
single-factor theorem (§4.3)\strut
\end{minipage} & \begin{minipage}[t]{0.38\columnwidth}\raggedright
\nolinkurl{step61_mode_t_single_factor_clean_separation_theorem_*}\strut
\end{minipage} & \begin{minipage}[t]{0.34\columnwidth}\raggedright
symbolic 3-step proof; non-vacuous at \(N=4\)\strut
\end{minipage}\tabularnewline
\begin{minipage}[t]{0.21\columnwidth}\raggedright
\(N_{\mathrm{gen}}\)-blindness theorem (§4.6)\strut
\end{minipage} & \begin{minipage}[t]{0.38\columnwidth}\raggedright
\nolinkurl{step62_mode_t_ngen_blindness_theorem_*}\strut
\end{minipage} & \begin{minipage}[t]{0.34\columnwidth}\raggedright
all-\(N\) invariance; verified to \(N=1000\); odd-doublet control \(N\)-sensitive\strut
\end{minipage}\tabularnewline
\begin{minipage}[t]{0.21\columnwidth}\raggedright
unification cluster (§4.4)\strut
\end{minipage} & \begin{minipage}[t]{0.38\columnwidth}\raggedright
\nolinkurl{step63_*}\strut
\end{minipage} & \begin{minipage}[t]{0.34\columnwidth}\raggedright
\(\sin^2\theta_W = \operatorname{Tr}T_3^2/\operatorname{Tr}Q^2 = 2/(16/3) = 3/8\); \(k_Y=5/3\); product control \(3/23\)\strut
\end{minipage}\tabularnewline
\begin{minipage}[t]{0.21\columnwidth}\raggedright
fork resolution (§4.4)\strut
\end{minipage} & \begin{minipage}[t]{0.38\columnwidth}\raggedright
\nolinkurl{step65_mode_b_proton_monopole_fork_resolution_*}\strut
\end{minipage} & \begin{minipage}[t]{0.34\columnwidth}\raggedright
record-stable breaking count \(=0\); fork object verified = step-41 coset\strut
\end{minipage}\tabularnewline
\begin{minipage}[t]{0.21\columnwidth}\raggedright
measurement extension (§7.4)\strut
\end{minipage} & \begin{minipage}[t]{0.38\columnwidth}\raggedright
\nolinkurl{step66/67/68_*}\strut
\end{minipage} & \begin{minipage}[t]{0.34\columnwidth}\raggedright
strict-extension witness; \(D\)-vs-\(\Sigma_f\) discriminator\strut
\end{minipage}\tabularnewline
\begin{minipage}[t]{0.21\columnwidth}\raggedright
\textbf{Prediction 3} (§6.3)\strut
\end{minipage} & \begin{minipage}[t]{0.38\columnwidth}\raggedright
\nolinkurl{step69_mode_t_record_stability_baryon_no_monopole_falsifiable_prediction_*}\strut
\end{minipage} & \begin{minipage}[t]{0.34\columnwidth}\raggedright
the \(2\times2\) table \(\{24,0;292,11674\}\); anti-circularity witnesses \nolinkurl{carrier_00827}, \nolinkurl{carrier_00770}; F27/F48 \(=0\) on the clean sample\strut
\end{minipage}\tabularnewline
\bottomrule
\end{longtable}

\hypertarget{the-qmgr-track-common-refinement-key-artifacts}{\subsection{The QM--GR track (common refinement) --- key artifacts}\label{the-qmgr-track-common-refinement-key-artifacts}}

\begin{longtable}[]{@{}lll@{}}
\toprule
\begin{minipage}[b]{0.21\columnwidth}\raggedright
result (§)\strut
\end{minipage} & \begin{minipage}[b]{0.38\columnwidth}\raggedright
artifact dir (under \nolinkurl{thread_qm_gr/steps/})\strut
\end{minipage} & \begin{minipage}[b]{0.34\columnwidth}\raggedright
headline numbers the validator re-derives\strut
\end{minipage}\tabularnewline
\midrule
\endhead
\begin{minipage}[t]{0.21\columnwidth}\raggedright
co-sourcing carrier + dynamics (§5.3)\strut
\end{minipage} & \begin{minipage}[t]{0.38\columnwidth}\raggedright
\nolinkurl{step25/26_*}\strut
\end{minipage} & \begin{minipage}[t]{0.34\columnwidth}\raggedright
one \(\psi\): Born + \(T[\psi]\) residuals \(0\); \(V=V_0+\kappa T_{00}\) fixed point, residual \(3.49\times10^{-16}\) (\(\kappa=0.3\)); runaway at \(\kappa=30\)\strut
\end{minipage}\tabularnewline
\begin{minipage}[t]{0.21\columnwidth}\raggedright
program theorems (§5.2)\strut
\end{minipage} & \begin{minipage}[t]{0.38\columnwidth}\raggedright
\nolinkurl{step30–38} family\strut
\end{minipage} & \begin{minipage}[t]{0.34\columnwidth}\raggedright
\(T_{\mathrm{QG\text{-}NoGo}}\); \(T_{\mathrm{QGR\text{-}Unique}}\); carrier-independence; nested control flips\strut
\end{minipage}\tabularnewline
\begin{minipage}[t]{0.21\columnwidth}\raggedright
RT bound + enrichment (§5.4)\strut
\end{minipage} & \begin{minipage}[t]{0.38\columnwidth}\raggedright
\nolinkurl{step41/42_*} (QM--GR numbering)\strut
\end{minipage} & \begin{minipage}[t]{0.34\columnwidth}\raggedright
\(S(A)\le\mathrm{mincut}\); saturation \(0.729/0.904/0.941\) at \(D{=}2,3,4\) (seeds 101--505)\strut
\end{minipage}\tabularnewline
\begin{minipage}[t]{0.21\columnwidth}\raggedright
entropy cone / MMI (§5.4)\strut
\end{minipage} & \begin{minipage}[t]{0.38\columnwidth}\raggedright
\nolinkurl{step44_holographic_mmi_entropy_cone_*}\strut
\end{minipage} & \begin{minipage}[t]{0.34\columnwidth}\raggedright
holographic \(I_3\le0\); min-cut \(I_3=0\); GHZ \(I_3=\log 2\)\strut
\end{minipage}\tabularnewline
\begin{minipage}[t]{0.21\columnwidth}\raggedright
discrete Einstein condition (§5.4)\strut
\end{minipage} & \begin{minipage}[t]{0.38\columnwidth}\raggedright
\nolinkurl{step45_discrete_einstein_consistency_*}\strut
\end{minipage} & \begin{minipage}[t]{0.34\columnwidth}\raggedright
\(38\) regions / \(14\) edges; rank \(10\), cokernel \(28\); residuals \(9\times10^{-17}\) vs \(0.042\)\strut
\end{minipage}\tabularnewline
\begin{minipage}[t]{0.21\columnwidth}\raggedright
premise door-test + grounding (§5.3)\strut
\end{minipage} & \begin{minipage}[t]{0.38\columnwidth}\raggedright
\nolinkurl{step47_common_carrier_door_test_*}\strut
\end{minipage} & \begin{minipage}[t]{0.34\columnwidth}\raggedright
per-law form-not-existence; commutator \(0.7071\); non-co-sourcing \(0.3806\)\strut
\end{minipage}\tabularnewline
\begin{minipage}[t]{0.21\columnwidth}\raggedright
ladder vs fork (§5.2)\strut
\end{minipage} & \begin{minipage}[t]{0.38\columnwidth}\raggedright
\nolinkurl{step48_ladder_vs_fork_resolution_*}\strut
\end{minipage} & \begin{minipage}[t]{0.34\columnwidth}\raggedright
ladder defect \(8\) (incl.~\(A_{\mathrm{RT}}=d_0+2d_2\) variant); fork \(0/0\); nested \(0\); route mismatch \(0.5\)\strut
\end{minipage}\tabularnewline
\begin{minipage}[t]{0.21\columnwidth}\raggedright
one fiber-volume ledger (§5.4)\strut
\end{minipage} & \begin{minipage}[t]{0.38\columnwidth}\raggedright
\nolinkurl{step50_born_area_one_fiber_volume_ledger_*}\strut
\end{minipage} & \begin{minipage}[t]{0.34\columnwidth}\raggedright
\(\mathrm{area}=\log 3\) seed-invariant; \(S_{\mathrm{Born}}\) varies; ratios \(0.9986/0.9973/0.9995 < 1\) strictly\strut
\end{minipage}\tabularnewline
\begin{minipage}[t]{0.21\columnwidth}\raggedright
monogamy forcing (§5.4)\strut
\end{minipage} & \begin{minipage}[t]{0.38\columnwidth}\raggedright
\nolinkurl{step51_f51_joint_prediction_*}\strut
\end{minipage} & \begin{minipage}[t]{0.34\columnwidth}\raggedright
geometric-dual compatibility (I3-independent); with-compat \(\max I_3=-0.033\); GHZ rel-gap \(0.79\); broken-compat un-forces\strut
\end{minipage}\tabularnewline
\begin{minipage}[t]{0.21\columnwidth}\raggedright
Bell / F49 (§7.1)\strut
\end{minipage} & \begin{minipage}[t]{0.38\columnwidth}\raggedright
\nolinkurl{step52_f49_bell_nonlocal_residual_*}\strut
\end{minipage} & \begin{minipage}[t]{0.34\columnwidth}\raggedright
CHSH \(2\sqrt2\) computed; bound \(2\) enumerated (16 strategies); control LHV table error \(1.8\times10^{-15}\)\strut
\end{minipage}\tabularnewline
\begin{minipage}[t]{0.21\columnwidth}\raggedright
BH information / F34 (§7.2)\strut
\end{minipage} & \begin{minipage}[t]{0.38\columnwidth}\raggedright
\nolinkurl{step53_f34_information_loss_normal_form_*}\strut
\end{minipage} & \begin{minipage}[t]{0.34\columnwidth}\raggedright
rank \(27=\) internal dim; explicit-\(G\) residual \(3\times10^{-15}\); partial witness \(5.3\times10^{-16}\) / gap \(0.0097\)\strut
\end{minipage}\tabularnewline
\begin{minipage}[t]{0.21\columnwidth}\raggedright
\(\Lambda\) / F50 (§7.3)\strut
\end{minipage} & \begin{minipage}[t]{0.38\columnwidth}\raggedright
\nolinkurl{step54_f50_background_moduli_demarcation_*}\strut
\end{minipage} & \begin{minipage}[t]{0.34\columnwidth}\raggedright
background sweep \(15/18\) pass; \(\kappa\) boundary (\(\le10\) pass, \(\ge15\) fail)\strut
\end{minipage}\tabularnewline
\begin{minipage}[t]{0.21\columnwidth}\raggedright
\textbf{Prediction 1} (§6.1)\strut
\end{minipage} & \begin{minipage}[t]{0.38\columnwidth}\raggedright
\nolinkurl{step55_f51_entanglement_underdetermines_geometry_*}\strut
\end{minipage} & \begin{minipage}[t]{0.34\columnwidth}\raggedright
fingerprint \(11{,}049\) features; central-diff rank/nullity \(9/5\) (seeds 101/202/303); clean shadows \(2\)/seed; tree control nullity \(0\)\strut
\end{minipage}\tabularnewline
\begin{minipage}[t]{0.21\columnwidth}\raggedright
\textbf{Prediction 2} (§6.2)\strut
\end{minipage} & \begin{minipage}[t]{0.38\columnwidth}\raggedright
\nolinkurl{step56_gravitational_mediation_bmv_prediction_*}\strut
\end{minipage} & \begin{minipage}[t]{0.34\columnwidth}\raggedright
initial entanglement \(6.2\times10^{-18}\); fork channel \(\mathcal N=0\), CHSH \(0\), damping \(1.0\); coherent control \(\mathcal N=0.3536\), CHSH \(\sqrt6=2.4495\); restriction derived from \nolinkurl{step48_access_tuple}+\nolinkurl{T_QG_NoGo}; \nolinkurl{mean_field_channel_used=false}\strut
\end{minipage}\tabularnewline
\bottomrule
\end{longtable}
\end{landscape}

\hypertarget{re-run-recipe}{\subsection{Re-run recipe}\label{re-run-recipe}}

{\small
\begin{verbatim}
cd physics_atlas/thread_<track>/steps/<step_dir>
python3 run_stepNN.py --self     # validator: recomputes + audits this step
python3 run_stepNN.py --chain    # additionally re-validates frozen upstream steps
python3 <build_script>.py        # regenerates every artifact deterministically
\end{verbatim}
}

All validators exit non-zero on any mismatch. The two review packets discussed in Appendix D additionally ship
\nolinkurl{SHA256SUMS} manifests verified from clean extracts.

\hypertarget{independent-re-derivation}{\subsection{Independent re-derivation}\label{independent-re-derivation}}

Beyond validator passes, every headline number in §6 was reproduced by code written independently of the build scripts
(direct SVD of centrally-differenced Jacobians for Prediction 1; direct BMV channel calculation for Prediction 2; direct recount of the \(2\times2\) partition and the
single-conjunct witness sets for Prediction 3), with exact agreement. The public reproducibility record for these
checks is the diagnostic-artifact manifest \nolinkurl{physics_atlas/INDEPENDENT_REDERIVATIONS.md}; process logs are not
part of the public artifact surface.
 \hypertarget{adversarial-audit-trail-what-was-rejected-and-why}{\section{adversarial audit trail: what was rejected, and why}\label{adversarial-audit-trail-what-was-rejected-and-why}}

We publish the program's rejected constructions deliberately. A framework as expressive as this one can manufacture
results in either direction --- toward novelty or toward the safety of recovery --- and the only durable defense is a
record showing both failure modes being caught. Every entry below was detected by the audit gates of §2.4/§3 (or by
line-by-line review against them), rejected, and replaced by the accepted construction cited in the body. None of the
rejected versions appears in any result.

\begin{landscape}

\hypertarget{tautologies-the-target-equal-to-its-evidence}{\subsection{Tautologies (the target equal to its evidence)}\label{tautologies-the-target-equal-to-its-evidence}}

\begin{longtable}[]{@{}llll@{}}
\toprule
\begin{minipage}[b]{0.19\columnwidth}\raggedright
rejected construction\strut
\end{minipage} & \begin{minipage}[b]{0.26\columnwidth}\raggedright
the defect\strut
\end{minipage} & \begin{minipage}[b]{0.21\columnwidth}\raggedright
the tell\strut
\end{minipage} & \begin{minipage}[b]{0.26\columnwidth}\raggedright
the accepted replacement\strut
\end{minipage}\tabularnewline
\midrule
\endhead
\begin{minipage}[t]{0.19\columnwidth}\raggedright
area--entropy ``law'' v1 (QM--GR)\strut
\end{minipage} & \begin{minipage}[t]{0.26\columnwidth}\raggedright
\emph{area} defined as boundary cell-count and \(S\) as cells \(\times \ln 2\) --- equality by construction\strut
\end{minipage} & \begin{minipage}[t]{0.21\columnwidth}\raggedright
validator would have passed; caught by reading the build: no can-fail case possible\strut
\end{minipage} & \begin{minipage}[t]{0.26\columnwidth}\raggedright
saturable bound \(S(A)\le\mathrm{mincut}\) with required strict-gap controls (§5.4)\strut
\end{minipage}\tabularnewline
\begin{minipage}[t]{0.19\columnwidth}\raggedright
area--entropy ``law'' v2\strut
\end{minipage} & \begin{minipage}[t]{0.26\columnwidth}\raggedright
bulk dimensions set equal to state dimensions, forcing the match\strut
\end{minipage} & \begin{minipage}[t]{0.21\columnwidth}\raggedright
same\strut
\end{minipage} & \begin{minipage}[t]{0.26\columnwidth}\raggedright
same\strut
\end{minipage}\tabularnewline
\begin{minipage}[t]{0.19\columnwidth}\raggedright
Born--area ``one ledger'' v1\strut
\end{minipage} & \begin{minipage}[t]{0.26\columnwidth}\raggedright
\(S_{F39}\) computed from the Schmidt spectrum, \(H_{F23}\) from eigenvalues of \(\rho_A\) --- \emph{the same spectrum twice} (von Neumann \(=\) Shannon of squared singular values)\strut
\end{minipage} & \begin{minipage}[t]{0.21\columnwidth}\raggedright
machine-\(\epsilon\) ``coincidence'' (\(8.9\times10^{-16}\)) --- an identity, true for any pure state, zero co-sourcing content\strut
\end{minipage} & \begin{minipage}[t]{0.26\columnwidth}\raggedright
geometric min-cut as the \(F39\) side: seed-invariant (\(\log 3\)) while \(S_{\mathrm{Born}}\) varies; strict saturation gaps; validator now \textbf{fails on} machine-equality (§5.4)\strut
\end{minipage}\tabularnewline
\begin{minipage}[t]{0.19\columnwidth}\raggedright
first-law ``consequence'' (δS relabeling)\strut
\end{minipage} & \begin{minipage}[t]{0.26\columnwidth}\raggedright
\(\delta S = \delta\langle H_{\mathrm{mod}}\rangle\) restated on a density matrix disconnected from the carrier --- no new checkable content\strut
\end{minipage} & \begin{minipage}[t]{0.21\columnwidth}\raggedright
correct but empty: nothing could have failed\strut
\end{minipage} & \begin{minipage}[t]{0.26\columnwidth}\raggedright
the entropy-cone/MMI consequence with the GHZ violator (§5.4)\strut
\end{minipage}\tabularnewline
\bottomrule
\end{longtable}

\hypertarget{circularity-and-hardcoding-the-load-bearing-clause-doing-no-work}{\subsection{Circularity and hardcoding (the load-bearing clause doing no work)}\label{circularity-and-hardcoding-the-load-bearing-clause-doing-no-work}}

\begin{longtable}[]{@{}llll@{}}
\toprule
\begin{minipage}[b]{0.19\columnwidth}\raggedright
rejected construction\strut
\end{minipage} & \begin{minipage}[b]{0.26\columnwidth}\raggedright
the defect\strut
\end{minipage} & \begin{minipage}[b]{0.21\columnwidth}\raggedright
the tell\strut
\end{minipage} & \begin{minipage}[b]{0.26\columnwidth}\raggedright
the accepted replacement\strut
\end{minipage}\tabularnewline
\midrule
\endhead
\begin{minipage}[t]{0.19\columnwidth}\raggedright
monogamy forcing v1 (F51)\strut
\end{minipage} & \begin{minipage}[t]{0.26\columnwidth}\raggedright
``forced\(\,=\,\)compatible \(\wedge\, I_3\le0\)'' with the compatibility flag \textbf{constantly true}; broken-compatibility control written as \textbf{literal} booleans\strut
\end{minipage} & \begin{minipage}[t]{0.21\columnwidth}\raggedright
the forcing reduced to the known fact ``holographic states obey MMI''; the control never computed\strut
\end{minipage} & \begin{minipage}[t]{0.26\columnwidth}\raggedright
compatibility as a computed, \(I_3\)-independent geometric-dual vector match; control computed over two constructed admissible sets --- the violator demonstrably survives without compatibility (§5.4)\strut
\end{minipage}\tabularnewline
\begin{minipage}[t]{0.19\columnwidth}\raggedright
clean-separation ``grounding'' v1 (SM)\strut
\end{minipage} & \begin{minipage}[t]{0.26\columnwidth}\raggedright
the internal grounding assumed clean separation\strut
\end{minipage} & \begin{minipage}[t]{0.21\columnwidth}\raggedright
anti-circularity gate: a hypothesis not independently satisfiable by a violating structure\strut
\end{minipage} & \begin{minipage}[t]{0.26\columnwidth}\raggedright
grounding one layer up in record-stability, with computed single-conjunct violators (\(60\), \(311\)) (§4.5)\strut
\end{minipage}\tabularnewline
\begin{minipage}[t]{0.19\columnwidth}\raggedright
SM ``candidate-law generation'' (early)\strut
\end{minipage} & \begin{minipage}[t]{0.26\columnwidth}\raggedright
a \texttt{total\_slots\ =\ 5} prior and an exterior/GUT-shaped predicate smuggled the answer\strut
\end{minipage} & \begin{minipage}[t]{0.21\columnwidth}\raggedright
caught by adversarial review + internal re-audit; demoted\strut
\end{minipage} & \begin{minipage}[t]{0.26\columnwidth}\raggedright
the token-blind neutral carrier and the conditional-selection chain (§4.2--4.3)\strut
\end{minipage}\tabularnewline
\begin{minipage}[t]{0.19\columnwidth}\raggedright
assorted (SM)\strut
\end{minipage} & \begin{minipage}[t]{0.26\columnwidth}\raggedright
a shape-flavored \nolinkurl{factor_local_breaking} predicate; a minimality ``currency'' rescue; an \(SU(2)\)-cubic-anomaly bug\strut
\end{minipage} & \begin{minipage}[t]{0.21\columnwidth}\raggedright
each detected in review; each would have faked a landing\strut
\end{minipage} & \begin{minipage}[t]{0.26\columnwidth}\raggedright
corrected chain (the published exclusions survive the fixes)\strut
\end{minipage}\tabularnewline
\bottomrule
\end{longtable}
\end{landscape}

\hypertarget{rigging-toward-a-desired-verdict-both-directions-prediction-1s-four-attempts}{\subsection{Rigging toward a desired verdict --- both directions (Prediction 1's four attempts)}\label{rigging-toward-a-desired-verdict-both-directions-prediction-1s-four-attempts}}

The entanglement-kernel prediction (§6.1) required four constructions; the first three were rejected.

\begin{enumerate}
\def\labelenumi{\arabic{enumi}.}
\tightlist
\item
  \textbf{Rigged toward the positive.} A bespoke hand-built graph (edge roles literally named \texttt{active\_bottleneck},
  \texttt{hidden\_half}) with three parallel channels --- the ``invisible mode'' engineered into the carrier; and the ``bulk
  invariant'' that moved was a single edge's weight relabeled as a distance. \emph{Defect class: manufactured witness.}
\item
  \textbf{Rigged toward the deflationary.} On the honest carrier, the genuine-mode dimension was \textbf{hardcoded to zero}
  (a dead conditional over an empty basis), and the scoring filter (``only multi-edge geodesics count; volume is
  slack'') was chosen so the computed kernel could not count. \emph{Defect class: verdict-by-filter; hardcoded score.}
\item
  \textbf{Inflated by a degeneracy artifact.} At the uniform-weight base point (all 14 edges equal, all cuts tied), a
  one-sided forward-difference Jacobian reported nullity 6 --- but every ``kernel'' direction changed the fingerprint at
  first order (residuals scaling linearly with step: \(2.3\times10^{-1}, 2.3\times10^{-3}, 2.3\times10^{-5}\) at
  \(\epsilon = 10^{-1,-3,-5}\)), and single edges were invisible in one direction only (cut slack). \emph{Defect class:
  linearization artifact at a non-generic point.}
\item
  \textbf{Accepted.} Generic tie-broken geometries; two-sided central differences; both-direction finite-range invisibility
  test; injective tree control; numbers reproduced independently of the build (nullity \(5/5/5\), shadows \(2/2/2\) on
  seeds 101/202/303).
\end{enumerate}

\hypertarget{what-the-trail-shows}{\subsection{What the trail shows}\label{what-the-trail-shows}}

Three observations. (i) The gates catch \emph{both} failure directions --- the same protocol rejected a manufactured ``new
physics'' witness and a manufactured ``it's just recovery'' verdict on the same question within successive attempts.
(ii) Validator-pass is necessary, never sufficient --- every catch above except the adversarial-review ones came from
reading the construction against the gates, then strengthening the validator to make the defect class mechanical.
(iii) The surviving §6 predictions are \emph{post-selection-hardened}: each now ships with a validator that fails on the
specific defect that almost produced it (machine-equality; one-sided Jacobians; degenerate base points; hardcoded
verdicts; filter-based scoring).
 \hypertarget{adversarial-review-record}{\section{adversarial review record}\label{adversarial-review-record}}

Two QM--GR construction results were stress-tested in adversarial review packets (cover letters
posing the sharpest objections we could state against ourselves; full artifacts; \texttt{SHA256SUMS}; offline re-run
instructions), and both were \textbf{settled}. The SM construction core separately passed an adversarial review cycle. These
cycles are part of the audit record; they are not journal peer review or independent replication. This
appendix records the questions asked, the verdicts, and the one methodological exchange of independent interest.

\hypertarget{bundle-1-the-ledgershadow-price-recognition-of-ryutakayanagi-5.4}{\subsection{Bundle 1 --- the ledger/shadow-price recognition of Ryu--Takayanagi (§5.4)}\label{bundle-1-the-ledgershadow-price-recognition-of-ryutakayanagi-5.4}}

\textbf{Question put to the review pass.} Is \texttt{area\ =\ shadow-price(entanglement)} a genuine, non-circular recognition landing ---
or a relabeling? Adversarial checklist included: is the saturable bound an identity in disguise; is the geometric side
secretly a function of the state; is the claimed subsumption of RT under the currency/constraint law fair to the
source material?

\textbf{Disposition.} Settled, two rounds. Round 1 returned required fixes (packaging defects; one scope-tightening of the
first-law step --- accepted and applied by reconstruction, not by prose). Round 2: \textbf{SETTLED --- genuine recognition
landing}; the review pass checked the framework source text and confirmed the subsumption claim is fair
(``saying max-flow/min-cut/bit-thread RT instantiates the currency-constraint/shadow-price schema is fair''), with the
stated non-claim boundary (recognizes/organizes RT; does not derive holography, \(1/4G\), Einstein equations, or
frame-transfer) endorsed as accurate.

\hypertarget{bundle-2-the-common-carrier-premise-and-the-fork-5.25.3}{\subsection{Bundle 2 --- the common-carrier premise and the fork (§5.2--5.3)}\label{bundle-2-the-common-carrier-premise-and-the-fork-5.25.3}}

\textbf{Question put to the review pass.} Do Steps 47--48 \emph{answer} the common-carrier objection or merely relocate it? Is the
co-sourcing warrant load-bearing or circular? Is the access structure of the ladder test fair to emergent-spacetime
positions, or baked in?

\textbf{Disposition.} Two rounds. Round 1: the fork graded \emph{genuine-conditional} (with a scope-naming request we adopted ---
the test covers weak-RT ladders, not strong reconstruction); the premise graded ``honest but insufficient --- it
relocates the objection.'' Round 2 (after a reconstruction pass that applied every precision fix while \textbf{declining} the
criterion shift --- see D.3): \textbf{SETTLED --- honest recognition source}; the review pass withdrew ``insufficient''
when the result is scored by the framework's own landing taxonomy, writing that demanding an internal derivation of an
inter-layer premise ``would indeed apply the wrong criterion,'' and confirmed: no number moved, the warrant extension is
genuine, conditionality consistently held.

\hypertarget{the-methodological-exchange-is-a-ground-landing-a-landing}{\subsection{The methodological exchange: is a GROUND landing a landing?}\label{the-methodological-exchange-is-a-ground-landing-a-landing}}

The substantive disagreement was about the \emph{grading criterion}, not the construction. The initial review position
--- a premise not derived is a premise not settled --- is the default of reductionist practice. The program's position,
which the review ultimately accepted: for an \emph{inter-layer} premise, the internal-derivation route is foreclosed in
principle (that is what makes it inter-layer); the honest landing mode is \textbf{GROUND} --- relocation up the tower to a
\emph{named, warranted, independently-checkable} source, carried with explicit source-warrant and theorem-strength caveats.
The criterion was applied consistently across both bundles (the same mode grounds the SM's clean-separation condition,
§4.5), and consistency --- not rhetoric --- is what settled the exchange. We record this because referees of this paper
will face the same choice of criterion, and the review record shows it can be resolved by inspection of the taxonomy
rather than by sympathy with the framework.

\hypertarget{the-sm-track-review}{\subsection{The SM track review}\label{the-sm-track-review}}

The SM selection-layer core (the neutral carrier, the token-blind selection, the closure chain) passed adversarial review
with verdict \emph{coherent finite construction} (sound-with-fixes; all fixes applied by reconstruction and re-validated).
The review confirmed: the can-fail controls have teeth, the adjudication is not label-rigged, and no hidden
cross-layer derivation is present. Notably, an early ``candidate-law generation'' claim was demoted \emph{during} this cycle
(a smuggled prior --- Appendix C.2), which we count as the review system working as intended.

\hypertarget{status-of-the-present-papers-predictions}{\subsection{Status of the present paper's predictions}\label{status-of-the-present-papers-predictions}}

The three §6 predictions postdate the reviewed bundles and have \textbf{not} yet received independent review; they have passed
the internal protocol of §3 (including independent re-derivation) and are published here precisely to invite that
review. The falsifiers in Table T1 are the review questions we would pose.
 \hypertarget{notation-the-six-primitives-at-a-glance}{\section{notation \& the six primitives at a glance}\label{notation-the-six-primitives-at-a-glance}}

\hypertarget{the-six-primitives}{\subsection{The six primitives}\label{the-six-primitives}}

\begin{longtable}[]{@{}llll@{}}
\toprule
\begin{minipage}[b]{0.11\columnwidth}\raggedright
label\strut
\end{minipage} & \begin{minipage}[b]{0.20\columnwidth}\raggedright
role\strut
\end{minipage} & \begin{minipage}[b]{0.37\columnwidth}\raggedright
operational test\strut
\end{minipage} & \begin{minipage}[b]{0.22\columnwidth}\raggedright
body appearances\strut
\end{minipage}\tabularnewline
\midrule
\endhead
\begin{minipage}[t]{0.11\columnwidth}\raggedright
\(P_1\)\strut
\end{minipage} & \begin{minipage}[t]{0.20\columnwidth}\raggedright
descent\strut
\end{minipage} & \begin{minipage}[t]{0.37\columnwidth}\raggedright
does \(t\) factor through \(q\)? (\(\mathcal O_q(t) = \varnothing\)?)\strut
\end{minipage} & \begin{minipage}[t]{0.22\columnwidth}\raggedright
blindness theorems (§4.6); baryon conservation (§6.3)\strut
\end{minipage}\tabularnewline
\begin{minipage}[t]{0.11\columnwidth}\raggedright
\(P_2\)\strut
\end{minipage} & \begin{minipage}[t]{0.20\columnwidth}\raggedright
representability\strut
\end{minipage} & \begin{minipage}[t]{0.37\columnwidth}\raggedright
what is admissibly expressible/selected at the layer\strut
\end{minipage} & \begin{minipage}[t]{0.22\columnwidth}\raggedright
the SM selection layer (§4.1)\strut
\end{minipage}\tabularnewline
\begin{minipage}[t]{0.11\columnwidth}\raggedright
\(P_3\)\strut
\end{minipage} & \begin{minipage}[t]{0.20\columnwidth}\raggedright
route mismatch\strut
\end{minipage} & \begin{minipage}[t]{0.37\columnwidth}\raggedright
\(E_1E_2 \ne E_2E_1\) on a shared carrier (\(\Delta_3^{\mathrm{comp}}\neq\varnothing\))\strut
\end{minipage} & \begin{minipage}[t]{0.22\columnwidth}\raggedright
quantize-vs-curve (§5.1)\strut
\end{minipage}\tabularnewline
\begin{minipage}[t]{0.11\columnwidth}\raggedright
\(P_4\)\strut
\end{minipage} & \begin{minipage}[t]{0.20\columnwidth}\raggedright
refinement\strut
\end{minipage} & \begin{minipage}[t]{0.37\columnwidth}\raggedright
a common refinement with commuting squares (F51)\strut
\end{minipage} & \begin{minipage}[t]{0.22\columnwidth}\raggedright
the parent \(L\) (§5.2); the GUT parent (§4.4)\strut
\end{minipage}\tabularnewline
\begin{minipage}[t]{0.11\columnwidth}\raggedright
\(P_5\)\strut
\end{minipage} & \begin{minipage}[t]{0.20\columnwidth}\raggedright
packaging\strut
\end{minipage} & \begin{minipage}[t]{0.37\columnwidth}\raggedright
closure of structure into an auditable unit\strut
\end{minipage} & \begin{minipage}[t]{0.22\columnwidth}\raggedright
closure chains (§4); record packets (§§6.2--6.3)\strut
\end{minipage}\tabularnewline
\begin{minipage}[t]{0.11\columnwidth}\raggedright
\(P_6\)\strut
\end{minipage} & \begin{minipage}[t]{0.20\columnwidth}\raggedright
audit\strut
\end{minipage} & \begin{minipage}[t]{0.37\columnwidth}\raggedright
the smuggle-check ledger\strut
\end{minipage} & \begin{minipage}[t]{0.22\columnwidth}\raggedright
the protocol (§3); every validator\strut
\end{minipage}\tabularnewline
\bottomrule
\end{longtable}

\emph{(Role labels only: no total algebra on \(\mathbb P\); no universal-decomposition claim --- both formal exclusions.)}

\hypertarget{recurring-symbols}{\subsection{Recurring symbols}\label{recurring-symbols}}

\begin{longtable}[]{@{}lll@{}}
\toprule
\begin{minipage}[b]{0.27\columnwidth}\raggedright
symbol\strut
\end{minipage} & \begin{minipage}[b]{0.49\columnwidth}\raggedright
meaning\strut
\end{minipage} & \begin{minipage}[b]{0.17\columnwidth}\raggedright
first use\strut
\end{minipage}\tabularnewline
\midrule
\endhead
\begin{minipage}[t]{0.27\columnwidth}\raggedright
\(H, q, Q, t\)\strut
\end{minipage} & \begin{minipage}[t]{0.49\columnwidth}\raggedright
carrier, access quotient, layer, readout\strut
\end{minipage} & \begin{minipage}[t]{0.17\columnwidth}\raggedright
§2.1\strut
\end{minipage}\tabularnewline
\begin{minipage}[t]{0.27\columnwidth}\raggedright
\(\mathcal O_q(t)\)\strut
\end{minipage} & \begin{minipage}[t]{0.49\columnwidth}\raggedright
descent obstruction set\strut
\end{minipage} & \begin{minipage}[t]{0.17\columnwidth}\raggedright
§2.1\strut
\end{minipage}\tabularnewline
\begin{minipage}[t]{0.27\columnwidth}\raggedright
\(\Delta_{\mathrm{fact}}(\pi_0,\pi_1)\)\strut
\end{minipage} & \begin{minipage}[t]{0.49\columnwidth}\raggedright
factorization defect (strict-extension witness set)\strut
\end{minipage} & \begin{minipage}[t]{0.17\columnwidth}\raggedright
§2.1\strut
\end{minipage}\tabularnewline
\begin{minipage}[t]{0.27\columnwidth}\raggedright
\(\Delta_3^{\mathrm{comp}}(E_1,E_2)\)\strut
\end{minipage} & \begin{minipage}[t]{0.49\columnwidth}\raggedright
route-mismatch defect of two idempotent completions\strut
\end{minipage} & \begin{minipage}[t]{0.17\columnwidth}\raggedright
§5.1\strut
\end{minipage}\tabularnewline
\begin{minipage}[t]{0.27\columnwidth}\raggedright
\(L = (d_0,d_1,d_2,d_3)\)\strut
\end{minipage} & \begin{minipage}[t]{0.49\columnwidth}\raggedright
the QM--GR common refinement; \(q_{\mathrm{QM}}=(d_0,d_1,d_2)\), \(q_{\mathrm{GR}}=(d_0,d_2,d_3)\)\strut
\end{minipage} & \begin{minipage}[t]{0.17\columnwidth}\raggedright
§5.2\strut
\end{minipage}\tabularnewline
\begin{minipage}[t]{0.27\columnwidth}\raggedright
\(\psi;\ \rho=|\psi|^2;\ T[\psi]\)\strut
\end{minipage} & \begin{minipage}[t]{0.49\columnwidth}\raggedright
co-sourcing field; Born readout; stress readout\strut
\end{minipage} & \begin{minipage}[t]{0.17\columnwidth}\raggedright
§5.3\strut
\end{minipage}\tabularnewline
\begin{minipage}[t]{0.27\columnwidth}\raggedright
\(V[\psi] = V_0 + \kappa T_{00}[\psi]\)\strut
\end{minipage} & \begin{minipage}[t]{0.49\columnwidth}\raggedright
semiclassical back-reaction (consistent at \(\kappa=0.3\); runaway at \(30\))\strut
\end{minipage} & \begin{minipage}[t]{0.17\columnwidth}\raggedright
§5.3\strut
\end{minipage}\tabularnewline
\begin{minipage}[t]{0.27\columnwidth}\raggedright
\(S(A)\), \(\mathrm{mincut}(A)\)\strut
\end{minipage} & \begin{minipage}[t]{0.49\columnwidth}\raggedright
entanglement entropy; geometric min-cut (``area'')\strut
\end{minipage} & \begin{minipage}[t]{0.17\columnwidth}\raggedright
§5.4\strut
\end{minipage}\tabularnewline
\begin{minipage}[t]{0.27\columnwidth}\raggedright
\(I_3(A{:}B{:}C)\)\strut
\end{minipage} & \begin{minipage}[t]{0.49\columnwidth}\raggedright
tripartite information; monogamy: \(\le 0\) holographically\strut
\end{minipage} & \begin{minipage}[t]{0.17\columnwidth}\raggedright
§5.4\strut
\end{minipage}\tabularnewline
\begin{minipage}[t]{0.27\columnwidth}\raggedright
\(\mathcal E(w)\), \(J\)\strut
\end{minipage} & \begin{minipage}[t]{0.49\columnwidth}\raggedright
entanglement fingerprint (\(11{,}049\) features); its Jacobian \(\partial\mathcal E/\partial w\)\strut
\end{minipage} & \begin{minipage}[t]{0.17\columnwidth}\raggedright
§6.1\strut
\end{minipage}\tabularnewline
\begin{minipage}[t]{0.27\columnwidth}\raggedright
\(\mathrm{RS}(C)\)\strut
\end{minipage} & \begin{minipage}[t]{0.49\columnwidth}\raggedright
record-stability (substrate \(\wedge\) capacity \(\wedge\) distinguishability)\strut
\end{minipage} & \begin{minipage}[t]{0.17\columnwidth}\raggedright
§6.3\strut
\end{minipage}\tabularnewline
\begin{minipage}[t]{0.27\columnwidth}\raggedright
F\(n\)\strut
\end{minipage} & \begin{minipage}[t]{0.49\columnwidth}\raggedright
layer-agnostic law \(n\) from the catalog (Appendix A.5)\strut
\end{minipage} & \begin{minipage}[t]{0.17\columnwidth}\raggedright
§2.2\strut
\end{minipage}\tabularnewline
\begin{minipage}[t]{0.27\columnwidth}\raggedright
\(T_{\mathrm{QG\text{-}NoGo}}\), \(T_{\mathrm{QGR\text{-}Unique}}\)\strut
\end{minipage} & \begin{minipage}[t]{0.49\columnwidth}\raggedright
program theorems: no fused/derived QG object; uniqueness of \(L\)\strut
\end{minipage} & \begin{minipage}[t]{0.17\columnwidth}\raggedright
§5.2\strut
\end{minipage}\tabularnewline
\bottomrule
\end{longtable}

\hypertarget{landing-mode-glossary-one-line-each}{\subsection{Landing-mode glossary (one line each)}\label{landing-mode-glossary-one-line-each}}

\textbf{COMPUTE} --- it drops out of the construction. \textbf{SELECT} --- chosen, on declared conditions. \textbf{PROVE-BLIND} --- proven
underivable ⇒ observed-input (a positive demarcation). \textbf{GROUND} --- supplied as the down-shadow of a named higher
source (conditional; caveats carried). \textbf{RECOGNITION} --- a known result reproduced as a law-instance, non-circularly.
\textbf{FORBIDDEN-RULE / DEMARCATION} --- a computed exclusion or law-fixable-vs-contingent classification, with a can-fail
control.

\hypertarget{the-headline-numbers-single-reference-row}{\subsection{The headline numbers (single reference row)}\label{the-headline-numbers-single-reference-row}}

\begin{longtable}[]{@{}ll@{}}
\toprule
\begin{minipage}[b]{0.47\columnwidth}\raggedright
quantity\strut
\end{minipage} & \begin{minipage}[b]{0.47\columnwidth}\raggedright
value\strut
\end{minipage}\tabularnewline
\midrule
\endhead
\begin{minipage}[t]{0.47\columnwidth}\raggedright
chiral gauge carrier; survivors\strut
\end{minipage} & \begin{minipage}[t]{0.47\columnwidth}\raggedright
\(11{,}990 \to \{2|3,\ SU(4)\}\) (\(\approx 99.3\%\) excluded)\strut
\end{minipage}\tabularnewline
\begin{minipage}[t]{0.47\columnwidth}\raggedright
embedding ratios (computed; control)\strut
\end{minipage} & \begin{minipage}[t]{0.47\columnwidth}\raggedright
\(\sin^2\theta_W = 3/8\), \(k_Y = 5/3\) (product control \(3/23\))\strut
\end{minipage}\tabularnewline
\begin{minipage}[t]{0.47\columnwidth}\raggedright
record-stable \(\times\) branch table\strut
\end{minipage} & \begin{minipage}[t]{0.47\columnwidth}\raggedright
\(\{24, \mathbf{0}; 292, 11{,}674\}\)\strut
\end{minipage}\tabularnewline
\begin{minipage}[t]{0.47\columnwidth}\raggedright
anti-circularity witnesses\strut
\end{minipage} & \begin{minipage}[t]{0.47\columnwidth}\raggedright
substrate-only \(60\); capacity-only \(311\)\strut
\end{minipage}\tabularnewline
\begin{minipage}[t]{0.47\columnwidth}\raggedright
ladder/fork defects\strut
\end{minipage} & \begin{minipage}[t]{0.47\columnwidth}\raggedright
ladder \(8\) (route mismatch \(0.5\)); fork \(0/0\); nested control \(0\)\strut
\end{minipage}\tabularnewline
\begin{minipage}[t]{0.47\columnwidth}\raggedright
premise controls\strut
\end{minipage} & \begin{minipage}[t]{0.47\columnwidth}\raggedright
commutator \(0.7071\); non-co-sourcing \(0.3806\); fixed point \(3.5\times10^{-16}\)\strut
\end{minipage}\tabularnewline
\begin{minipage}[t]{0.47\columnwidth}\raggedright
RT saturation trend (\(D=2,3,4\))\strut
\end{minipage} & \begin{minipage}[t]{0.47\columnwidth}\raggedright
\(0.729,\ 0.904,\ 0.941\)\strut
\end{minipage}\tabularnewline
\begin{minipage}[t]{0.47\columnwidth}\raggedright
MMI\strut
\end{minipage} & \begin{minipage}[t]{0.47\columnwidth}\raggedright
holographic \(I_3 \le 0\); GHZ \(I_3 = +\log 2\)\strut
\end{minipage}\tabularnewline
\begin{minipage}[t]{0.47\columnwidth}\raggedright
discrete Einstein condition\strut
\end{minipage} & \begin{minipage}[t]{0.47\columnwidth}\raggedright
\(38\) regions, \(14\) edges, rank \(10\), cokernel \(28\); residuals \(9\times10^{-17}\) vs \(0.042\)\strut
\end{minipage}\tabularnewline
\begin{minipage}[t]{0.47\columnwidth}\raggedright
CHSH\strut
\end{minipage} & \begin{minipage}[t]{0.47\columnwidth}\raggedright
quantum \(2\sqrt2\); enumerated local bound \(2\); classical control \(\sqrt2\)\strut
\end{minipage}\tabularnewline
\begin{minipage}[t]{0.47\columnwidth}\raggedright
BH-information classification\strut
\end{minipage} & \begin{minipage}[t]{0.47\columnwidth}\raggedright
rank \(27 =\) internal dim; explicit-\(G\) \(3\times10^{-15}\); partial-readout gap \(0.0097\)\strut
\end{minipage}\tabularnewline
\begin{minipage}[t]{0.47\columnwidth}\raggedright
\(\Lambda\) sweep\strut
\end{minipage} & \begin{minipage}[t]{0.47\columnwidth}\raggedright
\(15/18\) backgrounds pass; \(\kappa\) boundary between \(10\) and \(15\)\strut
\end{minipage}\tabularnewline
\begin{minipage}[t]{0.47\columnwidth}\raggedright
\textbf{Prediction 1 kernel}\strut
\end{minipage} & \begin{minipage}[t]{0.47\columnwidth}\raggedright
\(\operatorname{rank} J = 9\), \(\dim\ker J = 5\), clean shadows \(2\) (each of seeds 101/202/303); tree control \(0\)\strut
\end{minipage}\tabularnewline
\begin{minipage}[t]{0.47\columnwidth}\raggedright
\textbf{Prediction 2 (BMV)}\strut
\end{minipage} & \begin{minipage}[t]{0.47\columnwidth}\raggedright
fork channel \(\mathcal N=0\)/CHSH \(0\)/damping \(1.0\); coherent control \(\mathcal N=0.3536\)/CHSH \(\sqrt6\)/damping \(0.293\)\strut
\end{minipage}\tabularnewline
\begin{minipage}[t]{0.47\columnwidth}\raggedright
\textbf{Prediction 3 count}\strut
\end{minipage} & \begin{minipage}[t]{0.47\columnwidth}\raggedright
record-stable \(\wedge\) breaking \(= 0\) of \(11{,}990\)\strut
\end{minipage}\tabularnewline
\bottomrule
\end{longtable}
 
\bibliographystyle{unsrtnat}
\begin{thebibliography}{40}
\providecommand{\natexlab}[1]{#1}
\providecommand{\url}[1]{\texttt{#1}}
\expandafter\ifx\csname urlstyle\endcsname\relax
  \providecommand{\doi}[1]{doi: #1}\else
  \providecommand{\doi}{doi: \begingroup \urlstyle{rm}\Url}\fi

\bibitem[Weinberg(1967)]{Weinberg1967}
Steven Weinberg.
\newblock A model of leptons.
\newblock \emph{Phys. Rev. Lett.}, 19:\penalty0 1264--1266, 1967.
\newblock \doi{10.1103/PhysRevLett.19.1264}.

\bibitem[Georgi and Glashow(1974)]{GeorgiGlashow1974}
Howard Georgi and Sheldon~L. Glashow.
\newblock Unity of all elementary-particle forces.
\newblock \emph{Phys. Rev. Lett.}, 32:\penalty0 438--441, 1974.
\newblock \doi{10.1103/PhysRevLett.32.438}.

\bibitem['t~Hooft and Veltman(1974)]{tHooftVeltman1974}
Gerard 't~Hooft and Martinus J.~G. Veltman.
\newblock One-loop divergencies in the theory of gravitation.
\newblock \emph{Ann. Inst. H. Poincar\'e Phys. Theor. A}, 20\penalty0
  (1):\penalty0 69--94, 1974.
\newblock URL \url{https://www.numdam.org/item/AIHPA_1974__20_1_69_0/}.

\bibitem[Goroff and Sagnotti(1985)]{GoroffSagnotti1985}
Marc~H. Goroff and Augusto Sagnotti.
\newblock Quantum gravity at two loops.
\newblock \emph{Phys. Lett. B}, 160:\penalty0 81--86, 1985.
\newblock \doi{10.1016/0370-2693(85)91470-4}.

\bibitem[Ryu and Takayanagi(2006)]{RyuTakayanagi2006}
Shinsei Ryu and Tadashi Takayanagi.
\newblock Holographic derivation of entanglement entropy from {AdS/CFT}.
\newblock \emph{Phys. Rev. Lett.}, 96:\penalty0 181602, 2006.
\newblock \doi{10.1103/PhysRevLett.96.181602}.

\bibitem[Van~Raamsdonk(2010)]{VanRaamsdonk2010}
Mark Van~Raamsdonk.
\newblock Building up spacetime with quantum entanglement.
\newblock \emph{Gen. Relativ. Gravit.}, 42:\penalty0 2323--2329, 2010.
\newblock \doi{10.1007/s10714-010-1034-0}.

\bibitem[Maldacena and Susskind(2013)]{MaldacenaSusskind2013}
Juan Maldacena and Leonard Susskind.
\newblock Cool horizons for entangled black holes.
\newblock \emph{Fortschr. Phys.}, 61:\penalty0 781--811, 2013.
\newblock \doi{10.1002/prop.201300020}.

\bibitem[Bose et~al.(2017)Bose, Mazumdar, Morley, Ulbricht, Toro\v{s},
  Paternostro, Geraci, Barker, Kim, and Milburn]{Bose2017}
Sougato Bose, Anupam Mazumdar, Gavin~W. Morley, Hendrik Ulbricht, Marko
  Toro\v{s}, Mauro Paternostro, Andrew Geraci, Peter Barker, M.~S. Kim, and
  Gerard Milburn.
\newblock Spin entanglement witness for quantum gravity.
\newblock \emph{Phys. Rev. Lett.}, 119:\penalty0 240401, 2017.
\newblock \doi{10.1103/PhysRevLett.119.240401}.

\bibitem[Marletto and Vedral(2017)]{MarlettoVedral2017}
Chiara Marletto and Vlatko Vedral.
\newblock Gravitationally induced entanglement between two massive particles is
  sufficient evidence of quantum effects in gravity.
\newblock \emph{Phys. Rev. Lett.}, 119:\penalty0 240402, 2017.
\newblock \doi{10.1103/PhysRevLett.119.240402}.

\bibitem[Langacker(1981)]{Langacker1981}
Paul Langacker.
\newblock Grand unified theories and proton decay.
\newblock \emph{Phys. Rep.}, 72:\penalty0 185--385, 1981.
\newblock \doi{10.1016/0370-1573(81)90059-4}.

\bibitem[Tsiokos(2026{\natexlab{a}})]{Tsiokos_FoundationsII}
Ioannis Tsiokos.
\newblock Six birds foundations {II}: Admissibility, meta-theory, and the
  exact-six program, 2026{\natexlab{a}}.
\newblock Zenodo, \doi{10.5281/zenodo.19672278}.

\bibitem[Tsiokos(2026{\natexlab{b}})]{Tsiokos_FoundationsIII}
Ioannis Tsiokos.
\newblock Six birds foundations {III}: A finite audited interaction calculus
  for {SBT}, 2026{\natexlab{b}}.
\newblock Six Birds Theory foundations series.

\bibitem[Tsiokos(2026{\natexlab{c}})]{Tsiokos_FoundationsIV}
Ioannis Tsiokos.
\newblock Six birds foundations {IV}: A catalog of layer-agnostic structural
  laws, 2026{\natexlab{c}}.
\newblock Zenodo, \doi{10.5281/zenodo.20713187}.

\bibitem[Tsiokos(2026{\natexlab{d}})]{Tsiokos_FoEC}
Ioannis Tsiokos.
\newblock Six birds: Foundations of emergence calculus, 2026{\natexlab{d}}.
\newblock Zenodo, \doi{10.5281/zenodo.18365949}.

\bibitem[Tsiokos(2026{\natexlab{e}})]{Tsiokos_WhyMath}
Ioannis Tsiokos.
\newblock Why mathematics even works, 2026{\natexlab{e}}.
\newblock Zenodo, \doi{10.5281/zenodo.20712761}.

\bibitem[Anderson(1972)]{Anderson1972}
Philip~W. Anderson.
\newblock More is different.
\newblock \emph{Science}, 177:\penalty0 393--396, 1972.
\newblock \doi{10.1126/science.177.4047.393}.

\bibitem[Bouchiat et~al.(1972)Bouchiat, Iliopoulos, and Meyer]{Bouchiat1972}
C.~Bouchiat, J.~Iliopoulos, and Ph. Meyer.
\newblock An anomaly-free version of {Weinberg}'s model.
\newblock \emph{Phys. Lett. B}, 38:\penalty0 519--523, 1972.
\newblock \doi{10.1016/0370-2693(72)90532-1}.

\bibitem[Pati and Salam(1974)]{PatiSalam1974}
Jogesh~C. Pati and Abdus Salam.
\newblock Lepton number as the fourth ``color''.
\newblock \emph{Phys. Rev. D}, 10:\penalty0 275--289, 1974.
\newblock \doi{10.1103/PhysRevD.10.275}.

\bibitem[Fritzsch and Minkowski(1975)]{FritzschMinkowski1975}
Harald Fritzsch and Peter Minkowski.
\newblock Unified interactions of leptons and hadrons.
\newblock \emph{Ann. Phys.}, 93:\penalty0 193--266, 1975.
\newblock \doi{10.1016/0003-4916(75)90211-0}.

\bibitem[Dirac(1931)]{Dirac1931}
Paul A.~M. Dirac.
\newblock Quantised singularities in the electromagnetic field.
\newblock \emph{Proc. R. Soc. Lond. A}, 133:\penalty0 60--72, 1931.
\newblock \doi{10.1098/rspa.1931.0130}.

\bibitem['t~Hooft(1974)]{tHooft1974monopole}
Gerard 't~Hooft.
\newblock Magnetic monopoles in unified gauge theories.
\newblock \emph{Nucl. Phys. B}, 79:\penalty0 276--284, 1974.
\newblock \doi{10.1016/0550-3213(74)90486-6}.

\bibitem[{Super-Kamiokande Collaboration} et~al.(2020){Super-Kamiokande
  Collaboration}, Takenaka, et~al.]{SuperK2020}
{Super-Kamiokande Collaboration}, A.~Takenaka, et~al.
\newblock Search for proton decay via {$p \to e^+\pi^0$} and {$p \to
  \mu^+\pi^0$} with an enlarged fiducial volume in {Super-Kamiokande} {I--IV}.
\newblock \emph{Phys. Rev. D}, 102:\penalty0 112011, 2020.
\newblock \doi{10.1103/PhysRevD.102.112011}.

\bibitem[Witten(1982)]{Witten1982}
Edward Witten.
\newblock An {SU(2)} anomaly.
\newblock \emph{Phys. Lett. B}, 117\penalty0 (5):\penalty0 324--328, 1982.
\newblock \doi{10.1016/0370-2693(82)90728-6}.

\bibitem[Kobayashi and Maskawa(1973)]{KobayashiMaskawa1973}
Makoto Kobayashi and Toshihide Maskawa.
\newblock {CP}-violation in the renormalizable theory of weak interaction.
\newblock \emph{Prog. Theor. Phys.}, 49:\penalty0 652--657, 1973.
\newblock \doi{10.1143/PTP.49.652}.

\bibitem[Ford and Fulkerson(1956)]{FordFulkerson1956}
L.~R. Ford, Jr. and D.~R. Fulkerson.
\newblock Maximal flow through a network.
\newblock \emph{Can. J. Math.}, 8:\penalty0 399--404, 1956.
\newblock \doi{10.4153/CJM-1956-045-5}.

\bibitem[Freedman and Headrick(2017)]{FreedmanHeadrick2017}
Michael Freedman and Matthew Headrick.
\newblock Bit threads and holographic entanglement.
\newblock \emph{Commun. Math. Phys.}, 352:\penalty0 407--438, 2017.
\newblock \doi{10.1007/s00220-016-2796-3}.

\bibitem[Hayden et~al.(2013)Hayden, Headrick, and
  Maloney]{HaydenHeadrickMaloney2013}
Patrick Hayden, Matthew Headrick, and Alexander Maloney.
\newblock Holographic mutual information is monogamous.
\newblock \emph{Phys. Rev. D}, 87:\penalty0 046003, 2013.
\newblock \doi{10.1103/PhysRevD.87.046003}.

\bibitem[Greenberger et~al.(1989)Greenberger, Horne, and Zeilinger]{GHZ1989}
Daniel~M. Greenberger, Michael~A. Horne, and Anton Zeilinger.
\newblock Going beyond {Bell}'s theorem.
\newblock In Menas Kafatos, editor, \emph{Bell's Theorem, Quantum Theory, and
  Conceptions of the Universe}, pages 69--72. Kluwer Academic, Dordrecht, 1989.
\newblock \doi{10.1007/978-94-017-0849-4_10}.

\bibitem[Faulkner et~al.(2014)Faulkner, Guica, Hartman, Myers, and
  Van~Raamsdonk]{Faulkner2014}
Thomas Faulkner, Monica Guica, Thomas Hartman, Robert~C. Myers, and Mark
  Van~Raamsdonk.
\newblock Gravitation from entanglement in holographic {CFTs}.
\newblock \emph{J. High Energy Phys.}, 2014\penalty0 (3):\penalty0 51, 2014.
\newblock \doi{10.1007/JHEP03(2014)051}.

\bibitem[Maldacena(1998)]{Maldacena1998}
Juan~M. Maldacena.
\newblock The large {N} limit of superconformal field theories and
  supergravity.
\newblock \emph{Adv. Theor. Math. Phys.}, 2:\penalty0 231--252, 1998.
\newblock \doi{10.4310/ATMP.1998.v2.n2.a1}.

\bibitem[Bekenstein(1973)]{Bekenstein1973}
Jacob~D. Bekenstein.
\newblock Black holes and entropy.
\newblock \emph{Phys. Rev. D}, 7:\penalty0 2333--2346, 1973.
\newblock \doi{10.1103/PhysRevD.7.2333}.

\bibitem[Freivogel et~al.(2015)Freivogel, Jefferson, Kabir, Mosk, and
  Yang]{Freivogel2015}
Ben Freivogel, Robert~A. Jefferson, Laurens Kabir, Benjamin Mosk, and I-Sheng
  Yang.
\newblock Casting shadows on holographic reconstruction.
\newblock \emph{Phys. Rev. D}, 91:\penalty0 086013, 2015.
\newblock \doi{10.1103/PhysRevD.91.086013}.

\bibitem[Dyson(2013)]{Dyson2013}
Freeman Dyson.
\newblock Is a graviton detectable?
\newblock \emph{Int. J. Mod. Phys. A}, 28:\penalty0 1330041, 2013.
\newblock \doi{10.1142/S0217751X1330041X}.

\bibitem[Kafri et~al.(2014)Kafri, Taylor, and Milburn]{Kafri2014}
Dvir Kafri, Jacob~M. Taylor, and Gerard~J. Milburn.
\newblock A classical channel model for gravitational decoherence.
\newblock \emph{New J. Phys.}, 16:\penalty0 065020, 2014.
\newblock \doi{10.1088/1367-2630/16/6/065020}.

\bibitem[{Hyper-Kamiokande Collaboration}(2025)]{HyperK2025}
{Hyper-Kamiokande Collaboration}.
\newblock The {Hyper-Kamiokande} experiment: input to the update of the
  {European} strategy for particle physics, 2025.
\newblock URL \url{https://arxiv.org/abs/2506.16641}.

\bibitem[Bell(1964)]{Bell1964}
John~S. Bell.
\newblock On the {Einstein} {Podolsky} {Rosen} paradox.
\newblock \emph{Physics Physique Fizika}, 1:\penalty0 195--200, 1964.
\newblock \doi{10.1103/PhysicsPhysiqueFizika.1.195}.

\bibitem[Clauser et~al.(1969)Clauser, Horne, Shimony, and Holt]{CHSH1969}
John~F. Clauser, Michael~A. Horne, Abner Shimony, and Richard~A. Holt.
\newblock Proposed experiment to test local hidden-variable theories.
\newblock \emph{Phys. Rev. Lett.}, 23:\penalty0 880--884, 1969.
\newblock \doi{10.1103/PhysRevLett.23.880}.

\bibitem[Hawking(1976)]{Hawking1976}
Stephen~W. Hawking.
\newblock Breakdown of predictability in gravitational collapse.
\newblock \emph{Phys. Rev. D}, 14:\penalty0 2460--2473, 1976.
\newblock \doi{10.1103/PhysRevD.14.2460}.

\bibitem[Page(1993)]{Page1993}
Don~N. Page.
\newblock Information in black hole radiation.
\newblock \emph{Phys. Rev. Lett.}, 71:\penalty0 3743--3746, 1993.
\newblock \doi{10.1103/PhysRevLett.71.3743}.

\bibitem[Weinberg(1989)]{Weinberg1989}
Steven Weinberg.
\newblock The cosmological constant problem.
\newblock \emph{Rev. Mod. Phys.}, 61:\penalty0 1--23, 1989.
\newblock \doi{10.1103/RevModPhys.61.1}.

\end{thebibliography}
 

\end{document}
